\documentclass[11pt]{amsart}

\usepackage[brazil]{babel}
\usepackage{amsmath,amssymb,amsthm}
\usepackage[margin=1in]{geometry}
\usepackage[hyphens]{url}
\usepackage{hyperref}
\hypersetup{colorlinks=true,linkcolor=blue,citecolor=blue,urlcolor=blue}

\theoremstyle{plain}
\newtheorem{theorem}{Teorema}
\newtheorem{proposition}[theorem]{Proposi\c{c}\~ao}
\newtheorem{lemma}[theorem]{Lema}
\newtheorem{corollary}[theorem]{Corol\'ario}
\theoremstyle{definition}
\newtheorem{definition}[theorem]{Defini\c{c}\~ao}
\theoremstyle{remark}
\newtheorem{remark}[theorem]{Observa\c{c}\~ao}
\newtheorem{empresult}[theorem]{Resultado Emp\'irico}

\newcommand{\R}{\mathbb{R}}
\newcommand{\C}{\mathbb{C}}
\newcommand{\Z}{\mathbb{Z}}
\newcommand{\N}{\mathbb{N}}
\newcommand{\Q}{\mathbb{Q}}
\hyphenation{arXiv}

\title[A Conjectura da Cobertura Fraca de Wirsching]{A Conjectura da Cobertura Fraca de Wirsching: um
c\'alculo estendido, uma cota condicional e a estrutura do que resiste \`a cobertura}

\author{Renato Augusto Tavares}

\begin{document}

\maketitle

\vspace{-2.5em}
\begin{center}
\normalsize Universidade Federal de Goi\'as \par
\texttt{rat@discente.ufg.br} \par
ORCID: \url{https://orcid.org/0009-0002-0196-3311}
\end{center}

\begin{abstract}
A Conjectura da Cobertura Fraca de Wirsching, de 1998, pede um $K(l)$ subexponencial tal que um
conjunto $R_{j-1,j}\subset\Z_3^*$ suficientemente grande seja forçado a cobrir todo resíduo unidade
módulo $3^l$. Estendemos o cálculo exato de $j^*(l)$, o menor orçamento de cobertura, de $l=20$
para $l=23$, e comparamos quatro modelos de crescimento para $e(l):=j^*(l)-l\log_4 3$ contra sua
cauda. Um modelo constante é incondicionalmente impossível (provamos $e(l)\ge\tfrac14\log_2l-O(1)$)
e desfavorecido do início ao fim; as alternativas crescentes não são distinguidas entre si. Nossa
melhor cota, $j^*(l)\le\tfrac98l+\tfrac{33}2$, vem de um relaxamento por jogo de payoff médio via
uma correspondência, verificada computacionalmente, não provada em geral. A rota de Fourier é
bloqueada logo em seu primeiro passo: seu critério de positividade é comprovadamente inatingível em
todo orçamento; concedendo-o, cancelamento de raiz quadrada não conseguiria superar a cota do jogo
de payoff médio. Um cálculo adicional encontra a maior parte da massa das frequências primitivas em
coeficientes abaixo de um limiar fixo, evidência contra um reparo por conjunto excepcional esparso.
Diagnósticos de modelo nulo em um par encontram estrutura que nenhum dos nulos força, não
calibrada. Provamos que, para orçamentos $j>l$, o conjunto não coberto em um orçamento está contido
tanto em duas vezes quanto em quatro vezes o conjunto não coberto no orçamento anterior, dando uma
cota superior recursiva para $j^*(l)$. Essa cota é provada exata para $l+2\le j\le j^*(l)$ em
$l=3,\dots,13$ (verificado exaustivamente), empírica até $l=21$; o caso de fronteira $j=l+1$ é
provado em $l=3,\dots,6$ e empírico até $l=22$. Falsificamos uma regra de reparo de custo $1$ que
ele sugere, provamos que todo resíduo não coberto por último é $\equiv1\pmod3$ sempre que
$j^*(l)\ge l+1$ (todos os níveis calculados menos um), provamos uma lei de contenção de duas
classes mod $9$, e reportamos identidades adicionais, algumas verificadas mas não provadas. Duas
questões abertas são enunciadas com exatidão, mais estreitas que a conjectura; outras surgem em
outros pontos. A Conjectura da Cobertura Fraca permanece aberta.
\end{abstract}

\noindent\textbf{Palavras-chave:} mapa de Collatz, dinâmica $3$-ádica, conjectura de cobertura,
jogos de payoff médio, somas exponenciais, teoria de Littlewood--Offord.

\noindent\textbf{Classificação de Assuntos Matemáticos 2020:} 11B50, 11K70, 91A50, 11L07, 37A44.

\section{Introdução}
\label{sec:intro}

A monografia de 1998 de Wirsching~\cite{wirsching1998} estuda a densidade de conjuntos de
predecessores sob o mapa $3n+1$ através de uma construção de médias sobre os inteiros $3$-ádicos.
Uma cadeia de Markov limite construída a partir dessa construção tem uma densidade invariante, e a
própria rota de Wirsching até uma densidade positiva de predecessores passa por uma questão de
cobertura, enunciada nessa mesma monografia~\cite{wirsching1998}, sobre uma família de subconjuntos
finitos de $\Z_3^*$, as unidades $3$-ádicas. Um artigo posterior, do próprio autor, isola essa rota
como um alvo autocontido e enuncia, como conjecturas abertas, as lacunas ainda não transpostas entre
ela e uma prova~\cite{wirsching2003}. Para inteiros $j,k\ge0$, escreva
\begin{equation}
\label{eq:Rjk}
R_{j,k} := \Bigl\{\, \sum_{i=0}^{j} 2^{a_i}3^i \ \Bigm|\ j+k \ge a_0 > a_1 > \cdots > a_j \ge 0
\,\Bigr\} \subset \Z_3^*.
\end{equation}
Cada elemento de $R_{j,k}$ é construído a partir de uma sequência estritamente decrescente de $j+1$
``posições de dígito'' escolhidas em $\{0,\dots,j+k\}$: a $i$-ésima posição escolhida, contando do
topo, contribui $2^{a_i}$ escalado por $3^i$. A Conjectura da Cobertura Fraca de Wirsching
(Conjectura~3.9, Capítulo~V,~\cite{wirsching1998}) pede um $K(l)$ subexponencial, ou seja,
$\lim_{l\to\infty} K(l)e^{-\gamma l}=0$ para toda constante $\gamma>0$ (definição do próprio
Wirsching, mesmo capítulo), tal que $|R_{j-1,j}| \ge K(l)\cdot 3^l$ já force $R_{j-1,j}$ a cobrir
todo resíduo unidade módulo $3^l$; a cardinalidade do lado esquerdo é a da própria família de
tuplas, antes da redução, que pela bijeção da Proposição~\ref{prop:lowerbound} vale $\binom{2j}{j}$;
o tamanho (a priori muito menor) de sua imagem módulo $3^l$ é uma quantidade diferente. Escrevendo
$j^*(l)$ para o menor $j$ tal que $R_{j-1,j}$ cobre $(\Z/3^l\Z)^*$ exatamente (uma quantidade que
existe para todo $l$, incondicionalmente; Proposição~\ref{prop:exist} abaixo), a forma forçante da
conjectura e a taxa de crescimento de $j^*(l)$ coincidem uma vez que a cobertura, ao ser alcançada,
persista em todo orçamento maior: isso vale para $j\ge l$ (Teorema~\ref{thm:doubling},
Seção~\ref{sec:holdouts}), o que cobre todo nível da Tabela~\ref{tab:jstar}, mas não está
estabelecido abaixo de $l$, o regime que a taxa conjecturada acabaria por alcançar. Com essa
ressalva, a conjectura diz respeito à taxa de crescimento de
\begin{equation}
\label{eq:el}
e(l) := j^*(l) - l\log_4 3,
\end{equation}
o excesso do orçamento de cobertura verdadeiro sobre a taxa de ordem principal que um argumento de
cardinalidade sozinho prediz ($|R_{j-1,j}| = \binom{2j}{j} \sim 4^j/\sqrt{\pi j}$, logo igualar
$2\cdot3^{l-1}$ resíduos exige $j \sim l\log_4 3$ em ordem principal). No limiar,
$K(l) \asymp |R_{j^*(l)-1,j^*(l)}|/3^l \asymp 4^{e(l)}/\sqrt{j^*(l)}$; a correção
$\sqrt{j^*(l)}$ é polinomial em $l$ em todo nível de fato calculado, onde $j^*(l)=O(l)$
(Tabela~\ref{tab:jstar}), embora nenhuma cota polinomial incondicional sobre $j^*(l)$ seja provada
aqui (a cota incondicional da Proposição~\ref{prop:exist} é exponencial; a cota linear do
Corolário~\ref{cor:mpgbound} é condicional). Essa direção direta não precisa de ressalva de
persistência alguma: $j^*(l)$ é definido como um mínimo, a conjectura (para qualquer $K$
subexponencial válido) torna o primeiro orçamento $j$ com $\binom{2j}{j}\ge K(l)\cdot3^l$ um
orçamento de cobertura, e um mínimo é no máximo qualquer elemento de seu próprio conjunto
definidor, logo $j^*(l)\le l\log_43+o(l)$ segue diretamente; combinada com a cota inferior
incondicional da Proposição~\ref{prop:lowerbound}, a veracidade da conjectura força $e(l)$ a
crescer estritamente mais devagar que linearmente, sem ressalva alguma anexada. A ressalva acima
pertence à \emph{equivalência} enunciada no início deste parágrafo, não a essa implicação em uma só
direção: construir um $K$ válido \emph{a partir} da densidade limiar exige que a cobertura persista
em todo orçamento além de $j^*(l)$, não provado abaixo de $l$. A recíproca, que $e(l)=o(l)$
estabeleceria por sua vez a conjectura, precisa da mesma persistência, que o
Teorema~\ref{thm:doubling} só dá para $j\ge l$; abaixo de $l$, o regime que um $e(l)$ verdadeiramente
sublinear acabaria por alcançar, a recíproca é aberta.

Um manuscrito anterior, não publicado, do autor calculou $j^*(l)$ exatamente para $l=1,\dots,20$
por enumeração direta e reportou a tabela resultante sem uma análise da taxa de crescimento. Este
artigo faz três coisas. Primeiro, estende essa tabela até $l=23$ e dá uma comparação estatística
apropriada de modelos candidatos de crescimento para $e(l)$ contra a cauda da tabela estendida
(Seções~\ref{sec:computation} e~\ref{sec:growth}). Segundo, dá a melhor cota superior sobre
$j^*(l)$ que encontramos, via uma reformulação do orçamento de cobertura como um jogo de payoff
médio, condicional à correspondência verificada empiricamente entre os dois, enunciada no Resultado
Empírico~\ref{er:soundness} (Seção~\ref{sec:mpg}). Terceiro, e com maior extensão, reporta uma
tentativa sustentada de fechar a lacuna até a taxa conjecturada diretamente, primeiro pela versão
ingênua de cancelamento uniforme da abordagem natural de Fourier e, uma vez que essa versão se
mostra malsucedida por uma razão verificável, por um estudo direto dos resíduos que de fato resistem
à cobertura (Seções~\ref{sec:gapA}--\ref{sec:further}). Essa segunda tentativa não resolve a
conjectura. Ela produz um teorema sobre esses resíduos, uma falsificação da forma de custo $1$ de
uma regra de reparo que o teorema poderia sugerir, e duas questões abertas nomeadas, cada uma mais
estreita que a própria conjectura e merecedora de um ataque dedicado por conta própria.

\section{O cálculo}
\label{sec:computation}

A Tabela~\ref{tab:jstar} dá $j^*(l)$ e $e(l)$ para $l=1,\dots,23$. Os valores para $l\le20$
reproduzem exatamente a tabela do manuscrito anterior; foram recalculados aqui de forma
independente por uma reimplementação em Rust feita do zero da busca de cobertura (um programa
dinâmico de rotação de bitsets sobre resíduos módulo $3^l$), validada contra a implementação
original em Python em todo $l\le20$ antes de ser confiada para a extensão. Cada $j^*(l)$ reportado
é o menor orçamento no qual o bitset de cobertura do próprio programa dinâmico está cheio, de modo
que a busca certifica tanto a cobertura em $j^*(l)$ quanto sua ausência em $j^*(l)-1$ internamente,
por construção, em todo nível em que roda com $j^*(l)-1\ge l$ (excetuando vacuamente $l=1$, onde
$j^*(1)-1=0$ não é sequer um orçamento em que $R_{j-1,j}$ esteja definido); a concordância em
$l\le20$ contra uma implementação escrita de forma independente é a evidência deste artigo de que a
própria busca está correta. $l=21$ foi desde então conferido de forma cruzada da mesma maneira: uma
implementação separada em Python, feita do zero (bitsets de inteiros grandes nativos, seguindo a
recorrência do programa dinâmico em Rust como uma especificação em vez de copiar seu código),
reproduz de forma independente $j^*(21)=25$, confirmando tanto que $j=24$ falha em cobrir quanto que
$j=25$ cobre. $l=22,23$ são levados adiante apenas por inferência, sem uma conferência cruzada
separada nesses níveis; ambos são novos.

A busca acima só tenta orçamentos $j\ge l$ no nível $l$, de modo que, por si só, certifica $j^*(l)$
apenas contra candidatos menores dentro dessa faixa.

\begin{lemma}[Nenhum orçamento menor cobre]
\label{lem:nosmaller}
Para $l=2,\dots,23$, nenhum $j$ com $1\le j<l$ tem $R_{j-1,j}$ cobrindo $(\Z/3^l\Z)^*$, de modo que
$j^*(l)\ge l$ vale diretamente, não apenas em relação à faixa que a busca acima tenta.
\end{lemma}

\begin{proof}
Se $R_{j-1,j}$ cobre $(\Z/3^L\Z)^*$ para algum $L\ge j$, ele também cobre $(\Z/3^j\Z)^*$: a redução
módulo $3^j$ se fatora através da redução módulo $3^L$, e o homomorfismo de anéis
$\Z/3^L\Z\to\Z/3^j\Z$ se restringe a uma sobrejeção de grupos de unidades (toda unidade módulo $3^j$
se levanta a uma unidade módulo $3^L$), de modo que a imagem de um conjunto completo de unidades
módulo $3^L$ é um conjunto completo de unidades módulo $3^j$.

$R_{0,1}=\{1,2\}$ tem $2$ elementos, aquém de $|(\Z/3^l\Z)^*|=2\cdot3^{l-1}\ge6$ para $l\ge2$, logo
$j=1$ nunca cobre em $l\ge2$: $j^*(l)\ge2$ para $l\ge2$.

Fixe $l$ com $2\le l\le23$ e suponha, por contradição, que algum $j_0$ com $2\le j_0<l$ tenha
$R_{j_0-1,j_0}$ cobrindo $(\Z/3^l\Z)^*$. A redução acima, com $L=l$, dá $R_{j_0-1,j_0}$ cobrindo
também $(\Z/3^{j_0}\Z)^*$, isto é, $j^*(j_0)\le j_0$. Mas $j_0<l\le23$, logo $j_0$ é ele mesmo um
nível que a busca acima já rodou, começando do próprio orçamento $j_0$, e a busca certifica
diretamente que o orçamento $j_0$ falha em cobrir $(\Z/3^{j_0}\Z)^*$ ali (esse é o conteúdo bruto da
entrada $j^*(j_0)>j_0$ da Tabela~\ref{tab:jstar}, independente de qualquer orçamento menor também
falhar, o que não é suposto). Contradição. Logo nenhum $j_0<l$ cobre, dando $j^*(l)\ge l$.
\end{proof}

\begin{table}[h]
\centering
\begin{tabular}{c|c|c||c|c|c||c|c|c}
$l$ & $j^*(l)$ & $e(l)$ & $l$ & $j^*(l)$ & $e(l)$ & $l$ & $j^*(l)$ & $e(l)$ \\ \hline
1 & 1 & 0.208 & 9 & 13 & 5.868 & 17 & 21 & 7.528 \\
2 & 4 & 2.415 & 10 & 15 & 7.075 & 18 & 22 & 7.735 \\
3 & 6 & 3.623 & 11 & 16 & 7.283 & 19 & 23 & 7.943 \\
4 & 7 & 3.830 & 12 & 17 & 7.490 & 20 & 24 & 8.150 \\
5 & 9 & 5.038 & 13 & 18 & 7.698 & 21 & 25 & 8.358 \\
6 & 10 & 5.245 & 14 & 19 & 7.905 & 22 & 26 & 8.565 \\
7 & 11 & 5.453 & 15 & 20 & 8.113 & 23 & 27 & 8.773 \\
8 & 12 & 5.660 & 16 & 20 & 7.320 & & & \\
\end{tabular}
\caption{$j^*(l)$ e $e(l) = j^*(l) - l\log_4 3$, $l=1,\dots,23$.}
\label{tab:jstar}
\end{table}

O único platô na tabela, $j^*(15)=j^*(16)=20$, não é um erro: $l=15\to16$ é o único passo, entre os
$22$ de $l=1$ a $l=23$, em que o orçamento de cobertura não aumenta, e é exatamente o evento que a
contagem de platôs da Seção~\ref{sec:growth} registra. A diferença inteira $j^*(l)-l$, uma prima
mais grosseira de $e(l)$, é $0,2,3,3,4,4,4,4,4,5,5,5,5,5,5,4,4,4,4,4,4,4,4$ para $l=1,\dots,23$:
ela nunca cresce além de $5$ na faixa calculada, e nada provado aqui descarta que esse teto
persista; a cota logarítmica da Proposição~\ref{prop:lowerbound} sobre $e(l)$ não a toca, já que um
$j^*(l)-l$ limitado faz $e(l)=(j^*(l)-l)+l(1-\log_43)$ crescer linearmente, confortavelmente acima
de um piso logarítmico. Se o teto persistisse, $j^*(l)=l+O(1)$ tornaria $e(l)$ linear e, pela
identidade notada após \eqref{eq:el}, tornaria $K(l)$ exponencial e falsificaria a conjectura, a
contrapositiva da direção direta sem ressalva ali enunciada. Essa planicidade é um fato sobre a
faixa verificada, não uma cota provada em nenhum dos dois sentidos.

$l=21$ cabe na RAM física uma vez corrigido um passo de paralelização que desperdiçava memória na
reimplementação em Rust. $l=22$ precisa de aproximadamente $84\,\mathrm{GiB}$ de estado, acima dos
$62\,\mathrm{GiB}$ de RAM física da máquina, de modo que um arquivo de swap de $500\,\mathrm{GiB}$
foi adicionado à máquina especificamente para este cálculo; a execução levou aproximadamente
$10{,}749\,\mathrm{s}$ ($\sim\!3$ horas), com E/S de swap como o custo dominante, contra
$748\,\mathrm{s}$ para $l=21$ sem swap. $l=23$ exigiu aproximadamente $263\,\mathrm{GiB}$ de estado,
confortavelmente dentro do mesmo orçamento de swap, e $375{,}615\,\mathrm{s}$ ($\sim\!104$ horas).

\section{O crescimento de $e(l)$}
\label{sec:growth}

Comparamos quatro formas funcionais candidatas para o crescimento de $e(l)$ contra a cauda
$l=10,\dots,23$ ($n=14$), cada uma ajustada por mínimos quadrados ordinários: $e(l)=c$ (constante),
$e(l)=a+b\ln l$ (logarítmica), $e(l)=a+b\sqrt l$ (raiz quadrada), e $e(l)=a+bl$ (linear lenta). A
Tabela~\ref{tab:aic} dá as estatísticas de ajuste, usando as formas gaussiano-MQO
$\mathrm{AIC}=n\ln(\mathrm{RSS}/n)+2k$, $\mathrm{BIC}=n\ln(\mathrm{RSS}/n)+k\ln n$ (RSS a soma dos
quadrados dos resíduos, $k$ o número de coeficientes de regressão, $1$ para o modelo constante e
$2$ para os outros três), e o RMSE validado por leave-one-out; como apenas as diferenças
$\Delta\mathrm{AIC}$, $\Delta\mathrm{BIC}$ são usadas abaixo, a convenção para uma constante
aditiva compartilhada entre os modelos não as afeta. $e(l)$ é uma sequência exata e determinística,
sem modelo de amostragem por trás, de modo que estas são lidas de forma descritiva do início ao
fim, como escores de ajuste penalizados sobre a tabela observada, e não como evidência estatística
calibrada.

\begin{table}[h]
\centering
\begin{tabular}{l|c|c|c|c}
modelo & $k$ & $\Delta\mathrm{AIC}$ & $\Delta\mathrm{BIC}$ & RMSE LOOCV \\ \hline
constante & 1 & 16.091 & 15.452 & 0.5200 \\
logarítmico & 2 & 1.512 & 1.512 & 0.2991 \\
raiz quadrada & 2 & 0.722 & 0.722 & 0.2909 \\
linear lento & 2 & 0.000 & 0.000 & 0.2837
\end{tabular}
\caption{Comparação de modelos para $e(l)$, $l=10,\dots,23$. $\Delta$AIC/BIC relativos ao modelo de
melhor ajuste (linear lento).}
\label{tab:aic}
\end{table}

Um modelo constante é desfavorecido de forma acentuada e consistente: $\Delta\mathrm{AIC}=16.09$,
$\Delta\mathrm{BIC}=15.45$, ambos crescendo a cada $l$ adicionado desde $l=21$. Os três modelos de
crescimento não são distinguidos entre si pela regra prática convencional
$\Delta\mathrm{AIC}>2$, com o linear lento tomando uma dianteira estreita tanto em AIC quanto em
validação cruzada leave-one-out; os três regressores candidatos ($\log l$, $\sqrt l$, $l$) são eles
mesmos altamente correlacionados nessa faixa (correlação par a par $\ge0.993$), de modo que esse
ranqueamento reflete principalmente o quão fina é a margem atual. Pela identidade notada
após~\eqref{eq:el}, um $e(l)$ linear tornaria $K(l)$ exponencial, falsificando a conjectura
diretamente. Catorze pontos altamente correlacionados não conseguem distinguir essa possibilidade
de um crescimento logarítmico ou de raiz quadrada, de modo que essa comparação não resolve a
questão em nenhum dos dois sentidos; ela apenas diz quais taxas de crescimento permanecem vivas.

Nada desse aparato é de fato necessário para descartar um $e(l)$ constante: como $\log_43<1$,
arredondar uma quantidade constante mais linear $l\log_43+c$ ao inteiro mais próximo só pode mudar
por $0$ ou $1$ de um $l$ para o próximo, e no entanto $j^*(1)=1$ e $j^*(2)=4$, um incremento de
$3$. Uma origem com $e(l)$ constante é incompatível já com as duas primeiras entradas da tabela,
sem exigir nenhum teste estatístico.

Um fato combinatório direto descarta a mesma cauda de forma ainda mais acentuada, restrito a onde o
modelo ainda não está morto sobre a tabela completa: o modelo de arredondamento-constante por trás
dele prediz todo incremento em $\{0,1\}$, o que a tabela completa já refuta (o incremento de $3$
notado acima, e incrementos de $2$ em outros pontos da tabela, por exemplo de $j^*(2)=4$ para
$j^*(3)=6$). Restrito à cauda $l=10,\dots,23$ que a comparação de modelos acima usa ($13$
incrementos sobre os mesmos $14$ pontos, todos de fato em $\{0,1\}$, consistentes com a própria
predição do modelo ali), um modelo de $c$ fixo $J_l:=\mathrm{arredonda}(l\log_43+c)$ força a
mudança total sobre quaisquer $13$ passos consecutivos a
$\{\lfloor13\log_43\rfloor,\lceil13\log_43\rceil\}=\{10,11\}$, para todo $c$ (um fato padrão sobre
o arredondamento de uma progressão aritmética de inclinação irracional, verificado diretamente
sobre uma malha fina de valores de $c$). A mudança observada é $j^*(23)-j^*(10)=27-15=12$, fora
desse par para todo $c$: nenhum modelo de arredondamento-constante consegue reproduzir essa cauda
de forma alguma, uma conclusão que não precisa de probabilidade nenhuma. $e(l)$ é uma sequência
exata e determinística do início ao fim, de modo que nada neste parágrafo carrega uma interpretação
amostral. Entre os $13$ incrementos, exatamente um é zero (o platô $l=15\to16$ notado acima); sob o
modelo de arredondamento-constante, $e(l+1)-e(l)$ só pode ser $1-\log_43\approx0.2075$ ou
$-\log_43\approx-0.7925$. A contagem de platôs acima e a comparação por MQO acima são ambas funções
da mesma sequência de incrementos de $14$ pontos, de modo que não são conferências independentes,
embora a observação do platô incida especificamente sobre o modelo de arredondamento-constante, um
objeto diferente da regressão constante de intercepto livre $e(l)=c$ que a Tabela~\ref{tab:aic}
ajusta.

Um $e(l)$ constante é desfavorecido descritivamente por toda medida acima. Ele é também,
incondicionalmente, impossível.

\begin{proposition}[Ilimitação de $e(l)$]
\label{prop:lowerbound}
$e(l) \ge \tfrac14\log_2 l - O(1)$.
\end{proposition}

\begin{proof}
Tuplas de expoentes distintas dão somas distintas: como os $a_i$ são distintos, $a_{j-1}$ é o único
mínimo, de modo que $V=\sum_i 2^{a_i}3^i$ tem valuação $2$-ádica exatamente $a_{j-1}$ (todo outro
termo é divisível por uma potência de $2$ estritamente maior), o que recupera $a_{j-1}$ diretamente
a partir de $V$. O termo completo $2^{a_{j-1}}3^{j-1}$ é então subtraído, e o processo itera sobre
os $j-1$ termos restantes. Assim $R_{j-1,j}$, em bijeção com $j$-subconjuntos do conjunto de $2j$
elementos $\{0,\dots,2j-1\}$, tem cardinalidade exata $\binom{2j}{j}$ (essa cardinalidade é o
próprio Corolário~1.8 de Wirsching, Capítulo~V,~\cite{wirsching1998}; o argumento de injetividade
acima é elementar e já aparece, para o mesmo tipo de somas $2^a3^b$, em~\cite{tao2011blog}).
Cobrir $(\Z/3^l\Z)^*$, de tamanho $2\cdot3^{l-1}$, exige que o domínio seja ao menos tão grande
quanto o contradomínio: $\binom{2j}{j} \ge 2\cdot3^{l-1}$ é necessário em $j=j^*(l)$ em particular,
já que $R_{j^*(l)-1,j^*(l)}$ cobre por definição de $j^*(l)$. A cota explícita
$\binom{2n}{n} \le 4^n/\sqrt{3n+1}$, válida para todo $n\ge1$, transforma isso em
$4^j \ge 2\cdot3^{l-1}\sqrt{3j+1}$, logo
\[
j \ge l\log_4 3 - \log_4 3 + \log_4 2 + \tfrac12\log_4(3j+1).
\]
Descartando o último termo, não negativo, já dá $j \ge l\log_4 3 - O(1)$, logo
$3j+1 = \Omega(l)$, portanto $\log_4(3j+1) \ge \log_4 l - O(1) = \tfrac12\log_2 l - O(1)$.
Realimentando isso na desigualdade exibida dá a cota enunciada.
\end{proof}

\section{Uma cota superior por jogo de payoff médio}
\label{sec:mpg}

\subsection*{O jogo de precisão total}

Fixe $l\ge1$. Uma \emph{jogada de precisão total} a partir de uma unidade $z_0$ módulo $3^l$ é uma
sequência de $l$ movimentos: no passo $i=0,\dots,l-1$, a partir de uma unidade $z_i$ módulo
$3^{l-i}$, um movimento de custo $d_i\ge0$ é \emph{legal} quando $2^{d_i}z_i\equiv1\pmod3$, e
produz $z_{i+1}:=T_{d_i}(z_i):=(2^{d_i+1}z_i-2)/3$, que então precisa ser uma unidade módulo
$3^{l-i-1}$ (automático em $i=l-1$, onde $3^{l-i-1}=1$). Escreva $J:=\sum_{i=0}^{l-1}d_i$ para o
custo total.

\begin{lemma}[Fundamentação]
\label{lem:ground}
Para uma jogada de precisão total legal a partir de $z_0$ com custos $d_0,\dots,d_{l-1}$ e custo
total $J$, defina $a_{l-1}:=0$ e $a_{i-1}:=a_i+d_i+1$ para $i=l-1,\dots,1$ (logo
$a_0=J-d_0+l-1$). Então
\[
2^{J+l-1}z_0 \equiv \sum_{i=0}^{l-1} 2^{a_i}3^i \pmod{3^l}.
\]
O lado direito é, por construção, um elemento de $U(l,a_0)$ (Seção~\ref{sec:further}) com expoente
mais baixo $a_{l-1}=0$.
\end{lemma}

\begin{proof}
Por indução no comprimento $n$ da jogada ($n=1,\dots,l$), enunciado para uma jogada de comprimento
$n$ começando em uma unidade $z$ módulo $3^n$ (correspondendo à notação do lema em $n=l$,
$z=z_0$). Escreva $J_n:=\sum_{i=0}^{n-1}d_i$ e $a_{n-1}:=0$, $a_{i-1}:=a_i+d_i+1$ para
$i=n-1,\dots,1$, como acima.

Caso base $n=1$: o único movimento tem $2^{d_0}z\equiv1\pmod3$, que é exatamente
$2^{a_0+d_0}z\equiv2^{a_0}3^0\pmod3$ já que $a_0=0$ quando $n=1$.

Passo indutivo: suponha que a afirmação vale no comprimento $n-1$ para a subjogada com custos
$d_1,\dots,d_{n-1}$ começando em $z':=T_{d_0}(z)$, uma unidade módulo $3^{n-1}$; escrevendo
$a'_i:=a_{i+1}$ para essa subjogada, $a'_{n-2}=a_{n-1}=0$ corresponde ao seu caso base, e sua
recursão $a'_{i-1}=a'_i+d_{i+1}+1$ corresponde a $a_i=a_{i+1}+d_{i+1}+1$ (a recursão original no
índice $i+1$), de modo que a hipótese de indução dá
\[
2^{a_1+d_1}z' \equiv \sum_{i=0}^{n-2}2^{a_{i+1}}3^i \pmod{3^{n-1}}.
\]
Pela definição de $T_{d_0}$, $2^{d_0+1}z\equiv3z'+2\pmod{3^n}$. Multiplicando pela unidade
$2^{a_1+d_1}$ dá
\[
2^{a_1+d_1+d_0+1}z \equiv 3\cdot2^{a_1+d_1}z' + 2^{a_1+d_1+1} \pmod{3^n}.
\]
Como $a_0=a_1+d_1+1$ (a recursão em $i=1$), o lado esquerdo é $2^{a_0+d_0}z$ e o último termo é
$2^{a_0}$. Multiplicando a hipótese de indução por $3$, ela se torna, módulo $3^n$,
$3\cdot2^{a_1+d_1}z'\equiv\sum_{i=1}^{n-1}2^{a_i}3^i$, logo
\[
2^{a_0+d_0}z \equiv 2^{a_0} + \sum_{i=1}^{n-1}2^{a_i}3^i = \sum_{i=0}^{n-1}2^{a_i}3^i \pmod{3^n},
\]
a afirmação no comprimento $n$, usando $a_0+d_0=J_n+n-1$ (telescopando a recursão).
\end{proof}

\begin{empresult}[O pior caso do jogo é $j^*(l)$]
\label{er:soundness}
\[
j^*(l) = \max_{z_0}\ \min\{J : \text{$z_0$ tem uma jogada de precisão total legal de custo total
$J$}\},
\]
o máximo sobre unidades $z_0$ módulo $3^l$. Verificado contra a tabela exata conhecida de $j^*(l)$
para $l=1,\dots,12$, sem discrepância; que uma jogada de precisão total legal exista a partir de
toda unidade $z_0$ (o mínimo dentro da expressão é tomado sobre um conjunto não vazio) é parte do
que essa verificação confirma empiricamente. Uma prova geral precisaria adicionalmente que o
sucessor permaneça uma unidade no próximo módulo em cada passo, uma condição que a checagem de
legalidade mod $3$ sozinha não decide. O Lema~\ref{lem:ground} não estabelece isso por si só: o
giro $2^{J+l-1}$ no lema depende de $z_0$ através de seu próprio $J$ ótimo, de modo que não é uma
única reetiquetagem do grupo de unidades independente de $z_0$, e não temos uma prova geral da
identidade exibida.
\end{empresult}

\subsection*{O relaxamento de janela $k$}

Um movimento de custo $d\ge0$ em $z$ é \emph{legal} quando $2^dz\equiv1\pmod3$, como acima, e
\emph{seguro} na janela $k$ quando, para $m:=3^k$, cada um dos três resíduos $z,z+m,z+2m$ módulo
$3m$ mapeia sob $T_d$ de volta a uma unidade módulo $m$. Um jogador minimizador, vendo apenas
$z\bmod3^k$, escolhe um $d$ legal e seguro; um adversário então revela o próximo dígito $3$-ádico
oculto $\epsilon\in\{0,1,2\}$, forçando o estado verdadeiro a $T_d(z+3^k\epsilon)\bmod3^k$. Ambas
as exigências importam: legalidade sozinha pode enviar uma unidade a $0$ (tome $z=1$, $d=0$:
$T_0(1)=0$), e o filtro de segurança é exatamente o que mantém o passeio dentro de $\Z_3^*$
independentemente de qual dígito o adversário fornece. Legalidade depende de $d$ apenas através de
$d\bmod2$ (já que $2$ tem ordem $2$ módulo $3$), e $2$ é raiz primitiva módulo $3^{k+1}$ para todo
$k\ge0$ (ordem $2\cdot3^k$, por uma indução elementar de elevação do expoente), de modo que
$2^{d+1}\bmod3^{k+1}$, e portanto $T_d(z)\bmod3^k$ e a própria checagem de segurança, dependem de
$d$ apenas através de $d\bmod(2\cdot3^k)$: estados sucessores se repetem com esse período à medida
que $d$ cresce, e todo representante além do menor em sua classe residual tem custo estritamente
maior para o mesmo efeito, de modo que nunca é usado por uma política ótima. O conjunto de ações em
cada estado é, portanto, sem perda, o subconjunto de $\{0,\dots,2\cdot3^k-1\}$ que ali é legal e
seguro, que é finito. Que esse subconjunto seja não vazio em todo estado unidade foi verificado
computacionalmente em todo $k=3,\dots,14$ usado abaixo; uma prova geral fica em aberto aqui.

Esse jogo de visibilidade restrita é, por construção, um \emph{jogo de payoff médio} de dois
jogadores sobre os $2\cdot3^{k-1}$ estados unidade módulo $3^k$, com os agora finitos movimentos
legais e seguros como as arestas do minimizador, as três escolhas de $\epsilon$ do adversário como
as arestas do maximizador, e payoff a média de longo prazo dos custos de movimento $d$. A teoria
clássica de tais jogos, iniciada por Ehrenfeucht e Mycielski~\cite{ehrenfeuchtmycielski1979},
garante um valor e uma política posicional ótima em geral; o Teorema~\ref{thm:mpg} abaixo não
precisa dessa generalidade, apenas das doze políticas específicas cujos $(\rho_k,C_k)$ resultantes
são reportados na Tabela~\ref{tab:mpg}, cada uma conferida de forma independente conforme descrito
na Seção~\ref{sec:code}.

Uma política de janela $k$, rodada contra a sequência real de dígitos, dá uma jogada de precisão
total legal em qualquer $z$ módulo $3^l$: no passo $i$, o estado remanescente verdadeiro é $z_i$
módulo $3^{l-i}$, de modo que a política vê o verdadeiro $z_i\bmod3^k$ diretamente sempre que
$l-i\ge k$, isto é, em todo passo exceto os últimos $k-1$. Uma vez que $l-i<k$, seja porque $l<k$
globalmente ou porque uma jogada com $l\ge k$ alcançou seus próprios últimos $k-1$ passos, fixe um
levantamento do estado verdadeiro a um resíduo módulo $3^k$, no primeiro desses passos, preenchendo
os $l-i$ dígitos remanescentes do estado verdadeiro com $k-(l-i)$ dígitos adicionais arbitrários, e
rode o jogo de janela $k$ adiante a partir desse levantamento sobre os próprios dígitos
remanescentes da jogada (não um novo levantamento a cada passo). Dois fatos separados mantêm isso
consistente com o estado verdadeiro. Primeiro, segurança, por construção, mantém a imagem sob
$T_d$ de cada uma das três possíveis extensões de um dígito de um estado seguro como uma unidade
módulo $3^k$, de modo que a política de janela $k$ sempre tem um próximo estado bem definido
$s_{i+1}:=T_d(s_i+3^k\epsilon_i)\bmod3^k$ sobre o qual agir para qualquer dígito de extensão
$\epsilon_i\in\{0,1,2\}$, real ou de preenchimento, correspondendo ao sucessor geral $s'(\epsilon)$
usado na prova do teorema abaixo. Segundo, a partir da fórmula que define $T_d(z)=(2^{d+1}z-2)/3$,
$T_d(z)\bmod3^{k'}$ depende apenas de $z\bmod3^{k'+1}$ para todo $k'\ge0$; como $k>l-i-1$ ao longo
desse regime, o termo $3^k\epsilon_i$ nunca afeta nenhum dígito abaixo da posição $k$, em
particular nenhum dos $l-i$ dígitos de baixa ordem da entrada $s_i$, que pelo fato de precisão são
exatamente o que determina os $l-i-1$ dígitos de saída de baixa ordem que este passo acompanha, de
modo que qual $\epsilon_i$ é usado, real enquanto o estado verdadeiro ainda fornece um ou
arbitrário de preenchimento uma vez que não fornece mais, é irrelevante para eles. Consequentemente,
se o estado de janela $k$ $s_i$ concorda com o verdadeiro $z_i$ em seus $l-i$ dígitos mais baixos,
como acontece por construção no primeiro passo de preenchimento e por indução depois disso, então
$s_{i+1}$ concorda com o verdadeiro $z_{i+1}$ em seus $l-i-1$ dígitos mais baixos independentemente
de $\epsilon_i$, um dígito a menos de concordância garantida a cada passo, correspondendo
exatamente a quantos dígitos verdadeiros restam. A legalidade do movimento escolhido no passo $i$
depende apenas de $z_i\bmod3$, o dígito mais baixo, que está entre os dígitos verdadeiros e
concordantes em cada um desses passos já que $l-i\ge1$ ao longo de todos eles (os passos rodam
$i=0,\dots,l-1$); no estado terminal $z_l$, alcançado após o último passo, $l-i=0$ e a exigência de
unidade módulo $3^0=1$ é vazia. Assim o movimento escolhido permanece legal contra o estado
verdadeiro de menor precisão em todo passo, e a relação sucessora usada abaixo vale em todo passo,
incluindo esses últimos $k-1$, com $\epsilon_i$ tomado como o próximo dígito real sempre que o
estado verdadeiro ainda tem um e um dígito fixo arbitrário depois disso. Escrevendo $s_i$ para
qualquer que seja o estado de janela $k$ que $\sigma_k$ veja dessa forma no passo $i$ (o verdadeiro
$z_i\bmod3^k$, ou um levantamento fixo dele), $s_{i+1}$ é por construção o sucessor $s_i'(\epsilon_i)$
sob o movimento escolhido por $\sigma_k$, de modo que $s_0,\dots,s_l$ é exatamente uma jogada da
forma sobre a qual a prova abaixo soma, com custos de movimento iguais aos da jogada real (esse
regime é parte do que a checagem da política contra o programa dinâmico de cobertura em $k=4,7$
contra $l$ até $14$ já exercitou, com zero movimentos ilegais ou inseguros; toda jogada ali passa
por esse mesmo regime de últimos $k-1$ passos independentemente de quão grande $l$ seja).

\begin{theorem}[Cota do relaxamento de janela, condicional ao Resultado Empírico~\ref{er:soundness}]
\label{thm:mpg}
Para cada $k=3,\dots,14$, a política certificada por trás da Tabela~\ref{tab:mpg} (um movimento
legal e seguro em todo estado unidade módulo $3^k$, custo médio $\rho_k$) dá um $C_k$ racional
explícito com $j^*(l) \le \rho_k\cdot l + C_k$ para todo $l$.
\end{theorem}

\begin{proof}
Fixe $k$ na faixa. A política $\sigma_k$, o valor $\rho_k$, e um potencial
$h_k:(\Z/3^k\Z)^*\to\Q$ são calculados por um solucionador aninhado de melhoria de estratégia de
Howard que busca sobre o conjunto de ações finito estabelecido acima, e o movimento escolhido por
$\sigma_k$ é conferido, estado por estado e dígito por dígito, contra
\[
d_{\sigma_k}(s) + h_k(s'(\epsilon)) \le \rho_k + h_k(s)
\]
em todo estado $s$ e todo dígito $\epsilon$ do adversário, onde $s'(\epsilon)$ é o sucessor sob o
movimento escolhido por $\sigma_k$. Somando a desigualdade ao longo de qualquer jogada de $l$
passos $s_0,\dots,s_l$, ela telescopa para
\[
\sum_{i=0}^{l-1} d_{\sigma_k}(s_i) \le l\rho_k + h_k(s_0) - h_k(s_l) \le l\rho_k + \bigl(\max h_k -
\min h_k\bigr) =: l\rho_k + C_k,
\]
independentemente de quais dígitos o adversário de fato fornece, já que a desigualdade do potencial
vale para todo $\epsilon$ em todo passo. Assim $\sigma_k$, rodada como acima, dá a toda unidade $z$
módulo $3^l$ uma jogada de precisão total legal de custo total no máximo $l\rho_k+C_k$; o custo
mínimo sobre todas as jogadas legais a partir de $z$ é, portanto, também no máximo $l\rho_k+C_k$,
logo também o pior caso sobre $z$. Apenas a direção
$j^*(l)\le\max_{z_0}\min\{J\}$ da igualdade do Resultado Empírico~\ref{er:soundness} é usada aqui,
para identificar esse pior caso com $j^*(l)$; a direção reversa não desempenha papel algum nesta
prova.
\end{proof}

$\sigma_k$, $h_k$ e $\rho_k$ são calculados como racionais por um solucionador aninhado de melhoria
de estratégia de Howard, que busca ações $d$ até um teto fixo de $40$ (bem acima do maior $d$ que
qualquer política certificada abaixo de fato usa em qualquer $k=3,\dots,14$, que é $11$). A
desigualdade do potencial acima é conferida apenas para o movimento escolhido por $\sigma_k$ em
cada estado, contra os três dígitos do adversário, nunca contra o resto do conjunto de ações; a
legalidade e a segurança desse movimento escolhido são verificadas diretamente a partir das
definições completas do jogo acima ($2^dz\equiv1\pmod3$ e a condição dos três levantamentos), que
não fazem referência alguma ao teto, de modo que a cota superior resultante $j^*(l)\le\rho_kl+C_k$
também não depende dele. O solucionador também extrai uma cota inferior correspondente do
adversário a partir da mesma execução, mas esse cálculo abrange apenas o mesmo conjunto de ações
com teto: ``cota superior igual a cota inferior'' certifica $\rho_k$ exatamente como o valor do
jogo de payoff médio restrito a $d\le40$, não, sem argumento adicional, como o valor do jogo
irrestrito sobre o conjunto de ações completo $\{0,\dots,2\cdot3^k-1\}$ descrito acima, já que um
movimento com $d>40$, indisponível ao cálculo com teto da cota inferior, poderia em princípio
permitir ao jogador minimizador alcançar um valor verdadeiro menor do que o certificado com teto
descarta. A Tabela~\ref{tab:mpg} reporta $\rho_k$, $C_k$ sob esse entendimento: taxas certificadas
de uma política exibida e verificada, que é tudo o que o Teorema~\ref{thm:mpg} precisa, não um
valor provado-exato do jogo de janela $k$ completo. $\sigma_k$ e $\rho_k$ são validados de forma
cruzada por um solucionador independente de iteração de valor e, em $k=4$ e $k=7$, rodando a
política resultante contra o programa dinâmico de cobertura finito real para todo resíduo unidade
módulo $3^l$, $l$ até $14$, com zero movimentos ilegais ou inseguros encontrados.
A Tabela~\ref{tab:mpg} dá os valores calculados para $k=3,\dots,14$.

\begin{table}[h]
\centering
\begin{tabular}{c|c|c|c|c|c|c|c|c|c|c|c|c}
$k$ & 3 & 4 & 5 & 6 & 7 & 8 & 9 & 10 & 11 & 12 & 13 & 14 \\ \hline
$\rho_k$ & $2$ & $\tfrac53$ & $\tfrac32$ & $\tfrac32$ & $\tfrac75$ & $\tfrac{25}{19}$ & $\tfrac54$ &
$\tfrac{11}{9}$ & $\tfrac65$ & $\tfrac76$ & $\tfrac{119}{104}$ & $\tfrac98$ \\ \hline
$C_k$ & $5$ & $\tfrac{19}{3}$ & $\tfrac{15}{2}$ & $9$ & $10$ & $\tfrac{207}{19}$ & $\tfrac{47}{4}$ &
$\tfrac{115}{9}$ & $\tfrac{69}{5}$ & $\tfrac{44}{3}$ & $\tfrac{1621}{104}$ & $\tfrac{33}{2}$
\end{tabular}
\caption{A taxa certificada de relaxamento de janela $\rho_k$ e a constante $C_k$, $k=3,\dots,14$.}
\label{tab:mpg}
\end{table}

Em $k=14$, $\rho_{14} = 9/8 = 1.125$ e $C_{14}=33/2$, dando:

\begin{corollary}[Condicional ao Resultado Empírico~\ref{er:soundness}]
\label{cor:mpgbound}
$j^*(l) \le \tfrac98 l + \tfrac{33}2$.
\end{corollary}

Essa é, ao que sabemos, a melhor cota sobre a taxa de crescimento de $j^*(l)$ atualmente
disponível, condicional ao Resultado Empírico~\ref{er:soundness}. Observa-se que $\rho_k$ é não
crescente na faixa calculada $k=3,\dots,14$ (Tabela~\ref{tab:mpg}); não verificamos isso em geral,
e o Corolário~\ref{cor:mpgbound} usa apenas o único valor $k=14$. A técnica geral, representar um
problema combinatório do tipo cobertura como um jogo de payoff médio e calcular cotas racionais
exatas de relaxamento dessa forma, não é nova para esta aplicação: além do próprio teorema
fundacional de Ehrenfeucht e Mycielski~\cite{ehrenfeuchtmycielski1979}, um precedente metodológico
próximo e recente prova a racionalidade e computabilidade do raio de cobertura de um deslocamento
sófico primitivo por uma redução a jogo de payoff médio do mesmo tipo, aplicada a um objeto
diferente~\cite{meyerovitchyoung2026}; essa redução usa jogos de payoff médio não alternantes,
introduzidos pelos mesmos autores em~\cite{meyerovitchyoung2025nonalt}, uma variante dos jogos
alternantes clássicos de~\cite{ehrenfeuchtmycielski1979} usados aqui. Citamos os três como
precedente para o enquadramento; o que é novo aqui é sua aplicação específica a $j^*(l)$.

\begin{remark}
Ajustar $\rho_k = L + A/k$ por mínimos quadrados ordinários a $k=10,\dots,14$ dá $L\approx0.876$,
ainda bem acima da taxa verdadeira conjecturada $\log_4 3 = 0.7925\ldots$; o ajuste é sensível à
faixa usada (todo o intervalo $k=3,\dots,14$ dá $L\approx0.904$) e não é evidência de nenhum limite
específico. Apenas a direção $j^*(l)\le\rho_k l + C_k$, provada acima, está estabelecida. A
desigualdade reversa, $\inf_k \rho_k \le \limsup_l j^*(l)/l$, que seria necessária para concluir
qualquer coisa sobre para onde $\rho_k$ está caminhando, não está estabelecida.
\end{remark}

\section{A soma exponencial, e por que ela resiste a um ataque direto}
\label{sec:gapA}

A cota condicional do Corolário~\ref{cor:mpgbound}, $1.125\,l$, está bem acima da taxa conjecturada
$\log_4(3)\,l \approx 0.7925\,l$. A rota natural para fechar essa lacuna é analítico-Fourier. Aqui
$e(x):=e^{2\pi ix}$ denota o caractere aditivo padrão, sem relação com a sequência $e(l)$ de
\eqref{eq:el}. Por ortogonalidade em $\Z/3^l\Z$, o número de tuplas caindo na classe residual $z$ é
\begin{equation}
\label{eq:Nz}
N(z) = \frac1{3^l}\sum_{t=0}^{3^l-1} S(t)\,e\!\left(-\frac{tz}{3^l}\right), \qquad
S(t) := \sum_{j+k\ge a_0>\cdots>a_j\ge0} e\!\left(\frac{t}{3^l}\sum_{i=0}^j 2^{a_i}3^i\right),
\end{equation}
e $S(0)=|R_{j,k}|=:T$, já que $S(0)$ conta tuplas de expoentes e o argumento de injetividade na
prova da Proposição~\ref{prop:lowerbound} acima (tuplas distintas dão somas distintas, recuperando
$a_j$ a partir da valuação $2$-ádica da soma) se aplica sem alteração ao $(j,k)$ geral, não apenas
a $(j-1,j)$. A rota natural até a positividade é a desigualdade triangular: $N(z)>0$ para todo $z$
assim que $\sum_{t\ne0}|S(t)| < T$. Essa rota se fecha antes de se abrir. A
Proposição~\ref{prop:sqrtceiling} abaixo mostra que a soma à esquerda nunca fica abaixo de $T$, para
qualquer $l$, $j$, $k$: o critério não pode ser alcançado de forma alguma, para nenhum orçamento. O
restante desta seção estuda o que resta uma vez abandonada a rota de espectro completo: uma versão
de frequência primitiva da mesma ideia (\S5.2), e as razões adicionais abaixo pelas quais ela também
resiste.

\subsection{Um teto rígido sobre o cancelamento uniforme}

\begin{proposition}[O critério de espectro completo é inatingível]
\label{prop:sqrtceiling}
Para todo $l\ge1$ e todo $R_{j,k}$ com $T:=|R_{j,k}|>0$, $\sum_{t\ne0}|S(t)| \ge T$; em particular
$|S(3^{l-1})|\ge T/2$. Consequentemente o critério $\sum_{t\ne0}|S(t)|<T$ nunca é satisfeito, para
nenhum $l$, $j$, $k$, e tampouco a cota uniforme $|S(t)|=o(T/3^l)$ em todo $t$ não nulo, já que
isso forçaria $\sum_{t\ne0}|S(t)|=o(T)$.
\end{proposition}

\begin{proof}
Todo elemento de $R_{j,k}$ se reduz a $2^{a_0}\bmod3$ (todo outro termo carrega uma potência
positiva de $3$), logo é $\equiv1$ ou $2\pmod3$: nenhum elemento é divisível por $3$. Logo $N(z)=0$
para todo $z$ divisível por $3$, em particular $N(0)=0$, de modo que $\sum_{t=0}^{3^l-1}S(t)=3^lN(0)=0$
e $\sum_{t\ne0}S(t)=-S(0)=-T$; a desigualdade triangular dá
$\sum_{t\ne0}|S(t)|\ge\bigl|\sum_{t\ne0} S(t)\bigr|=T$.

O mesmo colapso se localiza: somando a fórmula de inversão sobre os $3^{l-1}$ resíduos divisíveis
por $3$, restam apenas as três frequências $t\in\{0,3^{l-1},2\cdot3^{l-1}\}$ na soma interna (todo
outro caractere é não constante nessa classe lateral do subgrupo de índice $3$ e tem média $0$),
dando $S(3^{l-1})+S(2\cdot3^{l-1})=-T$ exatamente, em todo $l$, $j$, $k$. Como
$S(3^l-t)=\overline{S(t)}$, os dois termos são conjugados, de modo que $\mathrm{Re}\,S(3^{l-1})=-T/2$
exatamente e $|S(3^{l-1})|\ge T/2$: um único par conjugado já força a cota, antes que o resto do
espectro seja somado.
\end{proof}

A Proposição~\ref{prop:sqrtceiling} fecha a rota ingênua diretamente, em todo $j$: não há limiar a
partir do qual ela passe a funcionar. Mesmo concedendo a premissa inatingível para efeito de
comparação, especializando $S(t)$, $T$ à família $R_{j-1,j}$ que a questão de cobertura deste
artigo trata ($T=\binom{2j}{j}$) e perguntando qual $j$ uma cota uniforme $|S(t)|\le C\sqrt{T}$
(cancelamento de raiz quadrada), aplicada via a desigualdade triangular sobre todas as $3^l-1$
frequências não nulas, forçaria dá $(3^l-1)C\sqrt{T} < T$, logo $T > ((3^l-1)C)^2$, isto é,
$4^j/\sqrt{3j+1} \ge T > ((3^l-1)C)^2$ (a cota explícita da prova da
Proposição~\ref{prop:lowerbound}), donde $j \ge 2\log_4(3)\,l - O(1) \approx 1.585\,l - O(1)$: já
\emph{pior} do que o $1.125\,l$ do Corolário~\ref{cor:mpgbound}, além de repousar sobre uma
premissa que a Proposição~\ref{prop:sqrtceiling} descarta em todo $j$. Ambos os fatos juntos
descartam a versão ingênua da abordagem de Fourier, e voltam a atenção para se um tratamento mais
fino, dependente de frequência (uma separação em arcos maiores e menores, na linguagem do método
do círculo), consegue fazer melhor.

\subsection{Um cálculo: a massa excepcional não permanece pequena}

Numericamente, $|S(t)|$ está longe de ser uniforme. Escrevendo $t=3^{l-c}t'$ com $3\nmid t'$
(valuação $l-c$), $t\cdot V/3^l=t'\cdot V/3^c$ para todo valor $V$ de tupla, de modo que $S(t)$ no
módulo $3^l$ é igual à mesma soma avaliada no módulo mais grosseiro $3^c$: uma frequência de
valuação $l-c$ é literalmente a soma de nível $c$ da mesma família de tuplas. Ao longo desta
subseção, $m$ denota um orçamento para a família $R_{m-1,m}$, o papel que $j$ desempenha no resto
deste artigo (não o parâmetro de profundidade $m:=3^k$ da Seção~\ref{sec:mpg}). A cota exata
$|S(3^{l-1})|\ge T/2$ da Proposição~\ref{prop:sqrtceiling} é apenas uma cota inferior, sem
reivindicação de maximalidade: em $m=1$, $l=2$ (a família $R_{0,1}=\{1,2\}$, $T=2$), $|S(3)|=1$
mas $|S(1)|=2\cos(\pi/9)\approx1.879$ excede isso, de modo que o par de topo nem sempre é a maior
frequência. Nos níveis diretamente verificados na família que esta seção de fato usa ($R_{m-1,m}$
no próprio limiar de cardinalidade de primeira ordem de cada nível: $l=10,m=9$; $l=12,m=11$;
$l=14,m=13$), o par de topo $t=3^{l-1},2\cdot3^{l-1}$ é observado como o maior, e valuação alta se
correlaciona com magnitude grande além dele, mas de forma frouxa: entre as doze maiores frequências
em $l=10$, as valuações são $9,9,8,8,7,7,8,8,8,8,7,7$, e em $l=12$ a maior frequência primitiva
(valuação $0$) já excede a mais fraca frequência de valuação $8$ verificada. O reparo natural é
restringir a frequências primitivas, excisando todo o espectro não primitivo (todo $c<l$, o par
acima incluído) e controlando o restante por uma cota mediada; se esse restante permanece esparso o
suficiente na escala que importa é uma questão que esta seção responde computacionalmente, para a
família $R_{j-1,j}$ em $l=18$, $j=16$ (o limiar de cardinalidade de primeira ordem, $T:=\binom{32}{16}$,
seis orçamentos abaixo do limiar de cobertura real $j^*(18)=22$: a cobertura necessariamente já
falha aqui, e este nível está abaixo da faixa $j\ge l$ em que as Seções~\ref{sec:holdouts}--\ref{sec:further}
trabalham por outro lado), com um resultado negativo. Escreva
$\|S\|_1:=\sum_{\gcd(t,3)=1}|S(t)|/T$ para a massa $L^1$ normalizada sobre frequências coprimas com
$3$. Essa restrição exclui $t=3^{l-1}$, cuja magnitude $|S(3^{l-1})|\sim T/\sqrt3$ é fixada pela
Proposição~\ref{prop:1547} e, empiricamente, é a maior nos três níveis verificados acima
($l=10,12,14$; não em geral, como o contraexemplo acima mostra), junto com sua conjugada
$t=2\cdot3^{l-1}$ (já que $S(3^l-t)=\overline{S(t)}$, essa frequência corresponde exatamente a ela
e é excluída pela mesma razão). A restrição a $\gcd(t,3)=1$ também exclui toda frequência não
primitiva, exatamente um terço do espectro, inteiramente, de modo que $\|S\|_1$ abaixo mede apenas
a massa das frequências primitivas. Aqui $\|S\|_1=5226.01$. Um limiar fixo faz os picos visíveis
parecerem esparsos, mas das $8{,}014$ frequências primitivas com $|S(t)|/T$ excedendo $0.001$, sua
contribuição combinada para $\|S\|_1$ é de cerca de $12.2$, de modo que mais de $99.7\%$ da massa
das frequências primitivas neste único nível calculado vive em uma população exponencialmente
numerosa de coeficientes em coeficientes abaixo de um limiar fixo, sem reivindicação alguma aqui
sobre como essa massa está distribuída dentro dela. Isso é evidência contra um reparo por conjunto
excepcional esparso entre as frequências primitivas neste nível acessível; nada diz sobre
frequências não primitivas, sobre todo $l$, e, já que $j=16$ está seis orçamentos abaixo de
$j^*(18)$, muito menos sobre a distribuição de massa perto do limiar de cobertura real.

Essa mesma estrutura do tipo $2^a3^b$ conecta o problema a uma discussão de 2011 de Terence Tao,
aplicando métodos de Littlewood--Offord e produtos de Riesz a um objeto construído a partir das
mesmas somas exponenciais do tipo $2^a3^b$, no contexto intimamente relacionado de um ciclo não
trivial hipotético de Collatz~\cite{tao2011blog}; sua construção é estruturalmente próxima de uma
fatia de fronteira de $R_{j,k}$ com o expoente mais baixo fixado em $0$. Sua própria conclusão, em
suas próprias palavras, é que ``uma vez que se usa a cota da teoria da transcendência, \dots{}
técnicas de Littlewood-Offord já não são de muita utilidade'', e que ele não conhece ``nenhuma
maneira de obter uma propriedade de separação não trivial entre potências de $2$ e potências de
$3$ além de via métodos da teoria da transcendência''.

\subsection{Um teste direto dos resíduos que de fato falham}

O teste mais direto examina os resíduos que de fato falham em ser cobertos, um passo de orçamento
antes de a cobertura ser bem-sucedida, isto é, $H(l,j^*(l)-1)$ (Seção~\ref{sec:holdouts}). Para
$l=10,\dots,15$ esse conjunto tem tamanho $2,1,3,1,3,1$ respectivamente, minúsculo frente a uma
média global de ocupação $\lambda=|R_{j^*(l)-2,j^*(l)-1}|/(2\cdot3^{l-1})$ que está ela mesma na
casa dos milhares nos mesmos níveis (de $1019$ em $l=10$ a $3695$ em $l=15$). Escreva a
\emph{célula de profundidade $c$} de um resíduo $z$ para sua classe módulo $3^c$, e sua
\emph{intensidade local de profundidade $c$} $\lambda_c(z)$ para a contagem total de acertos dentro
dessa classe dividida por $3^{l-c}$, o tamanho da classe (definida para $c<l$; em $c=l$ a classe é
$\{z\}$ ela mesma, e $\lambda_l(z)=0$ identicamente para qualquer retido, um caso degenerado
excluído abaixo). Um \emph{modelo de Poisson não homogêneo} ajustado a essa intensidade atribui a
$z$ a probabilidade de buraco $e^{-\lambda_c(z)}$. Calculado diretamente a partir do histograma
exato para $l=10,\dots,15$, cada um dos onze resíduos retidos ao longo desses seis níveis tem
$\lambda_c$ bem abaixo da média global em toda profundidade $c=8,9,10$ com $c<l$ (em $l=10$,
apenas $c=8,9$; $\lambda_c$ de $2.0$ a $108.4$, probabilidade de buraco de $e^{-2}$ a $e^{-108}$).
O modelo de Poisson não homogêneo atribui probabilidade $e^{-108}$ a um evento que ocorreu; $z$, os
níveis, e as profundidades verificadas são todos selecionados após ver os resíduos retidos, de modo
que este não é um teste pré-declarado e calibrado: o expoente é sugestivo, não uma rejeição formal,
e o modelo não é tratado como um candidato vivo abaixo com base nisso. A afirmação descritiva mais
fraca por trás disso continua de pé: resíduos retidos se situam sistematicamente em células de
baixa intensidade, com o quão abaixo da média variando amplamente por profundidade e nível. Na
profundidade mais fina que a definição permite, aquém do caso degenerado $c=l-1$ (uma classe de
apenas $3$ irmãos), o mesmo modelo usa a intensidade leave-one-out
$\lambda_{l-1}^{-z}(z):=\bigl(\text{soma das contagens de acertos dos outros $2$
irmãos}\bigr)/2$ no lugar de $\lambda_{l-1}(z)$, excluindo a própria contagem de $z$ e dividindo
pelo tamanho remanescente da classe (não $3$), para evitar o viés autorreferencial que uma classe
tão pequena carregaria de outra forma, para todo resíduo, não apenas para os retidos. Duas
conferências separadas seguem disso, tratando duas questões separadas. Primeiro, o modelo ainda
parece grosseiramente refutado nessa profundidade, do jeito que as profundidades mais grosseiras
acima pareciam (probabilidade de buraco até $e^{-108}$)? Não: a contagem total esperada de buracos
sob o modelo, somada sobre toda unidade, é $2.1$ contra $2$ resíduos retidos de fato observados em
$l=10$, $3.5$ contra $1$ em $l=11$, $4.5$ contra $3$ em $l=12$, $5.3$ contra $1$ em $l=13$, $5.0$
contra $3$ em $l=14$ ($l=15$ não rodado nessa profundidade); esse total agregado é invariante a
quais resíduos os buracos de fato são, de modo que não diz nada por si só sobre se a intensidade
acerta os resíduos corretos, apenas que o modelo não está mais grosseiramente descalibrado em
agregado nessa profundidade (os dois níveis em que o modelo mais superprediz, $l=11$ e $l=13$, são
eles mesmos evidência de que o ajuste é frouxo, não apertado). Segundo, a intensidade acerta quais
resíduos específicos resistem? Aqui a resposta é mais informativa: ranqueando toda unidade por essa
mesma intensidade leave-one-out (empates, onde a soma de valor inteiro coincide, quebrados de forma
arbitrária mas consistente, e irrelevantes para a conclusão já que nenhum resíduo retido está
dentro de um grande bloco empatado), os resíduos retidos reais se situam profundamente na cauda de
baixa intensidade em todo nível verificado dessa forma (ranques $12$ a $582$ entre $354{,}294$ a
$3{,}188{,}646$ unidades, apenas $l=12,13,14$), embora nem sempre entre os pouquíssimos resíduos de
ranque mais baixo. É essa conferência de ranque, e não o total agregado acima, que sustenta a
afirmação de correlação: lida nessa profundidade, a intensidade local se correlaciona fortemente
com quais resíduos resistem, uma questão que as profundidades mais grosseiras $c=8,9,10$ não
tratam por si sós, mas não a determina completamente. Um diagnóstico separado em um par menor e
mais tratável, $(l,m)=(14,16)$, para a família $R_{m-1,m}$ (contagem média de acertos $188.5$, três
orçamentos abaixo do limiar de cobertura $j^*(14)=19$), mantém todo $|S(t)|$ fixo e aleatoriza as
fases de pares conjugado-simétricos. Escrevendo $N_l(z)$ para a contagem exata de acertos do
resíduo $z$, $N_{l-1}$ para a contagem da mesma família de tuplas um nível abaixo, e
\[
F(z) := 3N_l(z)-N_{l-1}(z\bmod3^{l-1})
\]
para o desequilíbrio entre os três levantamentos de cada resíduo pai, restringir a frequências
primitivas $3\nmid t$ dá exatamente $\sum_{3\nmid t}S(t)e(-tz/3^l) = 3^{l-1}F(z)$, de modo que $F$
é identicamente $0$ em todo $z$ não unidade (tanto $N_l(z)$ quanto $N_{l-1}(z\bmod3^{l-1})$ se
anulam ali, já que nenhum elemento de uma família de tuplas é divisível por $3$). Sobre todo $z$, a
razão real entre o valor absoluto máximo dessa quantidade e sua raiz quadrada média é $15.79$,
enquanto oito embaralhamentos de fase independentes (semente fixa) dão razões em $[5.11,5.24]$ (um
segundo lote de oito, com semente independente, dá o intervalo mais largo $[4.98,5.40]$, de modo
que oito tentativas não fixam bem o nulo). Os embaralhamentos mantêm todo $|S(t)|$ fixo mas não
preservam a restrição de anulação-fora-de-unidades recém derivada, nem a integralidade, de modo que
não são, por si sós, o nulo correto.

A restrição de anulação-fora-de-unidades tem um significado algébrico direto para as fases.
Agrupando frequências primitivas em trios $t=r,\,r+3^{l-1},\,r+2\cdot3^{l-1}$ por resíduo pai $r$
módulo $3^{l-1}$ ($3\nmid r$), a anulação de $F$ no subgrupo $3(\Z/3^l\Z)$ é equivalente, pelo
mesmo colapso de soma de caracteres usado na Proposição~\ref{prop:sqrtceiling}, a
$S(t_1)+S(t_2)+S(t_3)=0$ para cada trio: três magnitudes fixas formando um triângulo fechado, não
degenerado exceto por uma fração desprezível (verificado diretamente em $(l,m)=(14,16)$: exatamente
dois dos $1{,}062{,}882$ trios primitivos de pai, um par conjugado, têm a maior magnitude igual à
soma das outras duas; a construção de rotação e reflexão do nulo abaixo só é possivelmente menos
aleatória nesses dois, sem efeito visível sobre uma estatística de $30$ tentativas). Isso fixa as
fases a uma família de um parâmetro, uma rotação comum da configuração real, junto com sua imagem
espelhada (negando as três fases antes de rotacionar, o que também soma zero). O parceiro conjugado
do trio no pai $r$ é o trio no pai $3^{l-1}-r$, e $S(3^l-t)=\overline{S(t)}$ (já usado na prova da
Proposição~\ref{prop:sqrtceiling}) liga os dois: apenas um trio por par conjugado recebe uma
rotação aleatória independente e um lançamento de moeda independente para a reflexão; as fases do
parceiro são então fixadas por conjugação, e não sorteadas por conta própria, de modo que o campo
resultante permanece real, correspondendo ao próprio $F$. Isso mantém todo $|S(t)|$ fixo e se anula
fora das unidades exatamente por construção, embora não a integralidade (o vetor real satisfaz
restrições adicionais que esse nulo não satisfaz, por exemplo $-N_{l-1}(r)\le F(z)\le2N_{l-1}(r)$
por pai, e ainda assim é mais extremo). Implementado e rodado em $(l,m)=(14,16)$ (30 tentativas,
registradas no repositório, Seção 9): o nulo restrito dá razões $\max/\mathrm{RMS}$ com média
$6.35$, faixa $[5.93,7.20]$, contra $15.79$ do vetor real, um fator residual de cerca de $2.5$. Por
Parseval, $\sum_z F(z)^2$ (logo a média quadrática sobre todo $z$) é a mesma para o vetor real e
para qualquer nulo de fase que preserve todo $|S(t)|$, restrito ou não, já que depende apenas das
magnitudes; um nulo que respeita a propriedade de anulação-fora-de-unidades coloca essa mesma massa
quadrada total apenas nas $2\cdot3^{l-1}$ posições unidade em vez de espalhá-la sobre todas as
$3^l$, o que é exatamente o que força $\mathrm{RMS}_{\mathrm{unidades}}=\sqrt{3/2}\cdot\mathrm{RMS}_{\mathrm{todas}}$
para tal nulo. ($\mathrm{RMS}_{\mathrm{unidades}}$ em si não entra em nenhum número reportado
abaixo; toda razão citada, real, irrestrita, ou restrita, usa $\mathrm{RMS}_{\mathrm{todas}}$ do
início ao fim, idêntica nas três por Parseval, de modo que toda a diferença entre $6.35$ e $5.2$
entre os nulos restrito e irrestrito vive no máximo, não em qual RMS é usado.) Essa identidade
governa apenas o RMS, não o máximo: Parseval nada diz sobre como concentrar essa mesma massa
quadrada em menos posições, mutuamente restritas, afeta uma estatística de valor extremo, e as
correlações de trio que a construção introduz poderiam em princípio empurrar $\max/\mathrm{RMS}$ em
qualquer direção. Que a razão medida do nulo restrito, $6.35$, caia perto de $\sqrt{3/2}\approx1.22$
vezes a faixa dos embaralhamentos irrestritos é uma propriedade empírica desta construção
específica, observada diretamente nas 30 tentativas, e não deduzida da identidade de Parseval
sozinha; o fator residual de $2.5$ repousa sobre esses números medidos, não sobre a identidade de
Parseval por si só. Tanto os lotes irrestritos de oito tentativas acima quanto este restrito de
trinta ainda são poucos demais para calibrar um nível de significância rigoroso (um teste baseado
em ranque contra trinta sorteios não consegue descer abaixo de $p\approx1/31$, por mais extremo que
seja o valor observado); o que o fator residual próximo de $2.5$ fornece, por si só, é evidência
contra as fases reais serem plenamente permutáveis com um nulo que respeita apenas $|S(t)|$ e a
restrição de anulação-fora-de-unidades, em vez de uma demonstração dedutiva (trinta sorteios de
fato poucos demais para calibrar um nível de significância, como acabamos de notar, podem sugerir
mas não estabelecer não permutabilidade). Isso não diz, por si só, que esse desvio precise de
estrutura de fase especificamente, já que nenhum dos dois nulos acima carrega informação alguma
sobre a intensidade local de nível $(l-1)$ documentada mais cedo nesta seção.

Um terceiro nulo testa diretamente a informação análoga de nível de pai, para a própria família
deste vetor: em vez de fixar $|S(t)|$, fixe o total bruto de cada pai, $N_{l-1}(r)$ para todo $r$, e
divida esse total multinomialmente $(\tfrac13,\tfrac13,\tfrac13)$ entre seus três levantamentos,
sem carregar informação alguma sobre as fases de frequência primitiva observadas (a divisão ainda
fixa os três levantamentos de cada pai a somarem seu total e não é uma afirmação sobre a intensidade
leave-one-out mais fina usada na conferência de ranque acima, uma propriedade de filhos individuais,
e não do total do pai que este nulo de fato condiciona). Rodado $30$ vezes (registrado no
repositório, Seção~9): esse nulo condicionado por total de pai dá razões $\max/\mathrm{RMS}$ com
média $17.4$, faixa $[14.9,21.1]$, excedendo o $15.79$ do vetor real em cerca de $83\%$ das
tentativas. A razão sozinha sugeriria que a aleatoriedade de total de pai reproduz a extremidade
observada; ela não sobrevive a um olhar sobre as duas quantidades absolutas que a razão divide.
Para cada levantamento $i$ de um total de pai $N$ sob esse nulo, $n_i$ é $\mathrm{Binomial}(N,\tfrac13)$,
de modo que o desequilíbrio $3n_i-N$ tem variância $2N$; somando os três levantamentos de todo pai
sobre o vetor inteiro dá $\mathbf{E}[\sum_z F(z)^2]=6T$ exatamente, $T=\binom{2m}{m}$. O
$\sum_z F(z)^2$ do vetor real é $9.35\times10^9$ contra uma expectativa de $6T=3.61\times10^9$, um
fator de $2.59$ e, contra a própria dispersão de tentativa a tentativa do nulo sobre as $30$
execuções (desvio padrão $5.6\times10^6$), cerca de $1{,}030$ desvios padrão do nulo: esse é o
desvio, um excesso de magnitude, o próprio $|S(t)|$ se desviando do que uma divisão independente
das contagens exatas de pai produz. O máximo absoluto do nulo, com média $477$ e nunca excedendo
$580$ sobre as $30$ tentativas, fica aquém do $698$ do vetor real em toda tentativa, mas isso não
acrescenta evidência adicional além do mesmo excesso: na própria raiz quadrada média do vetor real
($44.2$ contra $27.5$ do nulo, um fator de $1.61$), o máximo médio do nulo escalaria para cerca de
$477\times1.61\approx768$, acima do $698$ real, de modo que a comparação de máximo é o déficit de
raiz quadrada média reafirmado, não uma comparação independente. As próprias razões
$\max/\mathrm{RMS}$ concordam, ou melhor, o nulo excede a razão real na maioria das tentativas,
porque tanto o máximo quanto a raiz quadrada média ficam aquém por fatores comparáveis ($1.46$ e
$1.61$); nessa estatística livre de escala o nulo não é refutado de forma alguma, e a afirmação de
extremidade acima repousa inteiramente sobre o nulo de fase restrito, na mesma raiz quadrada média.
(Isso usa os próprios totais de pai do vetor $R_{15,16}$ diretamente, e não as figuras de
$\lambda_c$ reportadas mais cedo nesta seção, que vêm de uma família e profundidade diferentes,
$R_{j^*(l)-2,j^*(l)-1}$ em $c=8,9,10$; os dois são medições separadas do mesmo fenômeno
qualitativo, distorção grosseira nas contagens de tuplas, e não a mesma medição.) Assim o nulo
multinomial não remove a base para atribuir o excesso de energia observado a estrutura além dos
totais de pai; ele é esse excesso. Ele nada diz sobre fase, que é o próprio território do nulo
restrito acima: aquele nulo mantém todo $|S(t)|$ fixo por construção (Parseval então garante RMS
idêntico para o vetor real e o nulo igualmente) e ainda assim encontra o máximo real, $698$, bem
acima da média do nulo restrito, de cerca de $280$ nas mesmas unidades ($6.35$ vezes o RMS
compartilhado de $44.2$), um desvio puramente conduzido por fase que o nulo multinomial não pode
tratar já que nunca mantém $|S(t)|$ fixo para começar. Nenhum dos dois nulos, por si só ou juntos,
estabelece qual mecanismo, se algum, decide a cobertura, ou o quanto de cada desvio um teste
devidamente calibrado confirmaria. Nos níveis e no par verificados, a raridade de resíduos retidos
($l=10,\dots,15$) e a lacuna do embaralhamento de fase ($(l,m)=(14,16)$ apenas) não estão
estabelecidas como fenômenos independentes; ambas estão plausivelmente conectadas à mesma distorção
subjacente nas contagens de tuplas.

\begin{proposition}
\label{prop:1547}
Para todo $l\ge1$, $|S(3^{l-1})|/\binom{2m}{m} \to 1/\sqrt3$ quando $m\to\infty$, para a família
$R_{m-1,m}$.
\end{proposition}

\begin{proof}[Esboço da prova]
$\sum_i 2^{a_i}3^i \equiv 2^{a_0} \pmod 3$, de modo que $S(3^{l-1})/\binom{2m}{m}$ depende apenas
da distribuição limite da paridade do expoente de topo $a_0$ quando $m\to\infty$, que converge à
mistura $\tfrac13\omega + \tfrac23\omega^2$ das duas raízes cúbicas não triviais da unidade
$\omega=e^{2\pi i/3}$ (o intervalo limite $2m-1-a_0$ do máximo de um subconjunto aleatório de $m$
elementos de $\{0,\dots,2m-1\}$ é geométrico com parâmetro $1/2$, dando
$\Pr[a_0\text{ par}]\to1/3$). Essa mistura é igual a $-\tfrac12 - \tfrac{i\sqrt3}6$, de módulo
$1/\sqrt3$.
\end{proof}

$t/3^l=1/3$ exatamente, de modo que $l$ desaparece de $S(3^{l-1})$ por completo (apenas $V\bmod3$
importa), e o limite acima vale em toda escala $m$, incluindo ao longo da diagonal relevante para a
cobertura $m\sim l\log_4(3)$. Ele afia a cota exata $|S(3^{l-1})|\ge T/2$ da
Proposição~\ref{prop:sqrtceiling} até a razão assintótica exata $1/\sqrt3\approx0.577$, e confirma
que a massa que a Proposição~\ref{prop:sqrtceiling} força se concentra, nesta única frequência ao
menos, exatamente onde a numérica da \S5.2 a situa: na frequência de maior valuação, $t/3^l=1/3$.

\section{Um teorema sobre os resíduos que resistem à cobertura}
\label{sec:holdouts}

Escreva $H(l,j)$ para o conjunto das unidades módulo $3^l$ \emph{não} cobertas por $R_{j-1,j}$. Por
definição de $j^*(l)$ como o menor orçamento de cobertura, $H(l,j^*(l))=\emptyset$ e $H(l,j)\ne\emptyset$
para todo $j<j^*(l)$; se isso se estende a ``$j\ge j^*(l)$ exatamente quando $H(l,j)=\emptyset$''
exige que a vacuidade persista em todo orçamento maior também, o que não é automático apenas a
partir da definição. Todo resultado nesta seção precisa de $j\ge l$; $j^*(l) \ge l$ vale
diretamente para $l=1,\dots,23$ ($l=1$ pela Tabela~\ref{tab:jstar} diretamente, $l=2,\dots,23$ pelo
Lema~\ref{lem:nosmaller} acima; igualdade apenas em $l=1$), mas a taxa conjecturada $\log_4 3 < 1$
acabaria por colocar $j^*(l)$ abaixo de $l$, fora do alcance desta seção. Dentro dessa faixa, a
persistência de fato vale: o Teorema~\ref{thm:doubling} abaixo dá
$H(l,j)=\emptyset\Rightarrow H(l,j+1)\subseteq2\emptyset\cap4\emptyset=\emptyset$, de modo que
``$H(l,j)=\emptyset$ exatamente quando $j\ge j^*(l)$'' está estabelecido para todo $j\ge l$. Não
fazemos afirmação alguma sobre $j<l$.

\begin{theorem}[Inclusão de duplicação dos resíduos retidos]
\label{thm:doubling}
Para $j \ge l$, $H(l,j+1) \subseteq 2H(l,j) \cap 4H(l,j)$.
\end{theorem}

\begin{proof}
$R_{j-1,j}$ é o conjunto de valores $r=\sum_{i=0}^{j-1}2^{b_i}3^i$ sobre $j$-subconjuntos
$\{b_0>\cdots>b_{j-1}\}\subseteq\{0,\dots,2j-1\}$, e $R_{j,j+1}$ o mesmo com $(j{+}1)$-subconjuntos
de $\{0,\dots,2j+1\}$. Fixe $r\in R_{j-1,j}$ com conjunto de expoentes $S$, e $s\in\{1,2\}$.
Deslocando cada expoente para cima por $s$ dá $S+s\subseteq\{s,\dots,2j-1+s\}\subseteq\{1,\dots,2j+1\}$,
um $j$-subconjunto cujo valor induzido é $2^sr$ (um deslocamento uniforme de expoentes preserva o
ranque relativo, de modo que a potência de $3$ de cada termo permanece inalterada). Como $S+s$ não
contém $0$, adjuntar $0$ como o novo menor expoente dá um $(j{+}1)$-subconjunto de
$\{0,\dots,2j+1\}$: ele se torna o novo termo de ranque $j$, contribuindo $2^0\cdot3^j=3^j$. Assim
essa testemunha representa $2^sr+3^j$, e $3^j\equiv0\pmod{3^l}$ uma vez que $j\ge l$. Logo
$2R_{j-1,j}\cup4R_{j-1,j}\subseteq R_{j,j+1}\pmod{3^l}$ para $j\ge l$. Assim, se $z\in H(l,j+1)$,
isto é, $z$ não é representado no orçamento $j+1$, então $z$ não está nem em $2R_{j-1,j}$ nem em
$4R_{j-1,j}$: tanto $z/2$ quanto $z/4$ não são representados no orçamento $j$, isto é,
$z\in2H(l,j)\cap4H(l,j)$.
\end{proof}

\begin{corollary}[Contração de cadeia]
\label{cor:chain}
Para $j\ge l$: se $x, 2x, \dots, 2^{t-1}x \in H(l,j+1)$ são dois a dois distintos (todos reduzidos
módulo $3^l$, $t\ge1$), então $2^ix \in H(l,j)$ para todo $i=-2,\dots,t-2$ (isto é, $x/4,x/2$ e,
quando $t\ge2$, $x,2x,\dots,2^{t-2}x$), dois a dois distintos, uma cadeia de comprimento $t+1$.
\end{corollary}

\begin{proof}
Aplique o Teorema~\ref{thm:doubling} a cada um dos $t$ elementos: $2^ix\in H(l,j+1)$ dá
$2^{i-1}x, 2^{i-2}x\in H(l,j)$ para $i=0,\dots,t-1$. A união desses pares sobre todo $i$ é
exatamente $\{2^ix : i=-2,\dots,t-2\}$, $t+1$ elementos indexados por $t+1$ potências consecutivas
de $2$. Como $2$ tem ordem $2\cdot3^{l-1}$ módulo $3^l$ (um fato elementar de elevação do
expoente, já usado na Seção~\ref{sec:mpg}), quaisquer $2\cdot3^{l-1}$ potências consecutivas de $2$
vezes $x$ esgotam toda unidade módulo $3^l$, e qualquer sequência mais curta é dois a dois
distinta. A hipótese de que $x,2x,\dots,2^{t-1}x$ são dois a dois distintos já força
$t\le2\cdot3^{l-1}$, de modo que resta descartar a igualdade. Se $t=2\cdot3^{l-1}$, esses $t$
elementos seriam toda unidade módulo $3^l$, incluindo todo elemento da imagem (não vazia) de
$R_{j,j+1}$; mas $H(l,j+1)$, o complemento dessa imagem, não contém elemento algum assim. Logo
$t<2\cdot3^{l-1}$, e os $t+1$ elementos de saída, uma sequência de no máximo $2\cdot3^{l-1}$
potências consecutivas, são dois a dois distintos.
\end{proof}

\begin{corollary}[Bootstrap de comprimento de sequência]
\label{cor:bootstrap}
Escrevendo $\mathrm{maxrun}(H)$ para o maior $t$ tal que $H$ contém $t$ elementos dois a dois
distintos da forma $x,2x,\dots,2^{t-1}x$ (todos reduzidos módulo $3^l$), com
$\mathrm{maxrun}(\emptyset):=0$, $j^*(l) \le j + \mathrm{maxrun}(H(l,j))$ para todo $j\ge l$.
\end{corollary}

\begin{proof}
Seja $r:=\mathrm{maxrun}(H(l,j))$ e suponha, por contradição, $H(l,j+r)\ne\emptyset$; escolha
$z\in H(l,j+r)$, uma cadeia de comprimento $1$ ($t=1$) no sentido do Corolário~\ref{cor:chain}.
Aplicando o Corolário~\ref{cor:chain} uma vez em cada um dos $r$ orçamentos
$j+r-1,j+r-2,\dots,j$ (todos $\ge l$, já que $j\ge l$), uma cadeia de comprimento $t$ no orçamento
$j'+1$ se torna uma cadeia de comprimento $t+1$ no orçamento $j'$ a cada vez, de modo que, após $r$
aplicações, a cadeia de comprimento $1$ em $j+r$ se torna uma cadeia de comprimento $r+1$ em
$H(l,j)$, contradizendo $\mathrm{maxrun}(H(l,j))=r$. Logo $H(l,j+r)=\emptyset$, isto é,
$j^*(l)\le j+r=j+\mathrm{maxrun}(H(l,j))$.
\end{proof}

\begin{empresult}
\label{er:bootstrap}
O Corolário~\ref{cor:bootstrap} é empiricamente exato, $j^*(l) = j + \mathrm{maxrun}(H(l,j))$, em
todo $(l,j)$ com $l+1\le j\le j^*(l)$ que conseguimos calcular: $l=2$ verificado diretamente em seu
único orçamento não trivial, $j=l+1=3$ ($H(2,3)=\{7\}$ módulo $9$, $\mathrm{maxrun}=1$,
correspondendo a $j^*(2)=4$), $l=3,4$ verificados diretamente da mesma forma em $j=l+1$
($\mathrm{maxrun}(H(3,4))=2$ e $\mathrm{maxrun}(H(4,5))=2$, correspondendo a $j^*(3)=6$ e
$j^*(4)=7$, com o argumento de extensão de largura de fronteira seguindo a
Proposição~\ref{prop:tight} abaixo também provando essa mesma igualdade ali como um teorema, não
como uma verificação empírica), $l=5,\dots,21$ em todo orçamento assim, e $l=22$ no único orçamento
$j=l+1$, incluindo um caso com $|H(l,j)| > 3\times10^6$. (Além de $j^*(l)$, $H(l,j)=\emptyset$ e
$\mathrm{maxrun}(\emptyset)=0$, de modo que a igualdade exibida deixa de valer ali; a afirmação
diz respeito apenas a orçamentos até e incluindo $j^*(l)$.) O extremo inferior $j=l+1$ tem uma
justificativa real por trás: a igualdade no próprio $j=l$ vale para $l=1,\dots,5$ mas falha em todo
nível verificado, $l=6,\dots,13$ ($\mathrm{maxrun}(H(6,6))=5$, dando
$j+\mathrm{maxrun}=11\ne10=j^*(6)$; $\mathrm{maxrun}(H(7,7))=5$, dando $12\ne11=j^*(7)$;
$\mathrm{maxrun}(H(8,8))=8$, dando $16\ne12=j^*(8)$; $\mathrm{maxrun}(H(9,9))=8$, dando
$17\ne13=j^*(9)$; $\mathrm{maxrun}(H(10,10))=9$, dando $19\ne15=j^*(10)$;
$\mathrm{maxrun}(H(11,11))=8$, dando $19\ne16=j^*(11)$; $\mathrm{maxrun}(H(12,12))=8$, dando
$20\ne17=j^*(12)$; $\mathrm{maxrun}(H(13,13))=8$, dando $21\ne18=j^*(13)$); os
Corolários~\ref{cor:chain} e~\ref{cor:bootstrap} precisam ambos de $j\ge l$ para valer, mas isso
sozinho não explica por que $j=l$ especificamente começa a falhar em $l=6$. A subfaixa
$l+2\le j\le j^*(l)$ é provada em $l=3,\dots,13$: a Proposição~\ref{prop:tight} abaixo, que também
prova o orçamento de fronteira $j=l+1$ em $l=3,4,5,6$ especificamente. A partir de $l=7$, $j=l+1$
permanece apenas empírico, ligado à própria recíproca aberta da Proposição~\ref{prop:h014}.
\end{empresult}

Uma recíproca mais forte, por testemunha, do Teorema~\ref{thm:doubling}, de que $x/2,x/4\in H(l,j)$
já força $x\in H(l,j+1)$, o que por sua vez tornaria uma cadeia de duplicação máxima tanto
suficiente quanto necessária para a sobrevivência até um dado orçamento, é falsa como afirmação
irrestrita: $1547 \in 2H(7,8)\cap4H(7,8)$, e no entanto $1547 \notin H(7,9)$ (ambas as duas
testemunhas de orçamento $9$ de $1547$ têm uma configuração de expoentes que a construção de
deslocamento-e-adjunção do Teorema~\ref{thm:doubling} não reconstrói diretamente a partir de uma
testemunha de orçamento $8$). A Seção~\ref{sec:further} retoma uma forma restrita e mais fraca
dessa recíproca (redundância de canto, não a afirmação por testemunha que este contraexemplo
refuta), que permanece aberta em geral e é provada em níveis específicos.

Um fortalecimento natural do Teorema~\ref{thm:doubling}, de que toda testemunha individual admite
um reparo local de custo limitado, também é falso em sua forma mais otimista.

Chame de \emph{profundidade de Durfee} de uma testemunha $\{a_0>\cdots>a_{j-1}\}\in R_{j-1,j}$ no
orçamento $j$ a contagem $\#\{i : a_i \ge j\}$, o número de expoentes escolhidos no orçamento ou
acima dele.

\begin{empresult}[Falha da regra de reparo local de custo $1$]
\label{er:falsify}
Na transição de orçamento $(l,j)=(6,10)\to(7,11)$, chame uma unidade módulo $3^7$ de \emph{filha} e
sua redução módulo $3^6$ de sua \emph{mãe}; $100$ das $1458$ filhas (todas as unidades módulo
$3^7$) têm profundidade de Durfee mínima atingível excedendo a profundidade mínima de sua mãe em
mais de um, até três; em $(7,11)\to(8,12)$, $332$ das $4374$ filhas (todas as unidades módulo
$3^8$) fazem o mesmo, também até três. Uma verificação igualmente exaustiva define o \emph{custo de
reparo} de um resíduo filho $r$ com resíduo mãe $p$ como $\min_{S,S'}|S\triangle S'|$, a menor
diferença simétrica de conjuntos de expoentes sobre toda testemunha de conjunto de expoentes $S$
para $p$ no orçamento $j$ e toda testemunha de conjunto de expoentes $S'$ para $r$ no orçamento
$j+1$ ($|S|=j$, $|S'|=j+1$, de modo que essa distância é sempre ímpar, $\ge1$, e o mínimo abrange
toda a fibra de testemunhas em cada orçamento). Isso encontra custos de reparo até $7$ em ambas as
transições, contra um custo de $1$ para a esmagadora maioria. Como toda filha em ambas as
transições é verificada exaustivamente, isso falsifica um custo de reparo de exatamente $1$ nessas
duas transições diretamente; não estabelece que nenhum custo de reparo limitado exista.
\end{empresult}

Isso falsifica, em sua forma de custo $1$ originalmente proposta, uma conjectura natural que este
projeto vinha considerando anteriormente (de que a cobertura se propaga via uma regra uniforme de
reparo local de custo $1$, orçamento por orçamento); seu conteúdo válido remanescente é o
Teorema~\ref{thm:doubling} e seus corolários acima.

\section{Estrutura adicional dos últimos resíduos retidos}
\label{sec:further}

Os resíduos em $H(l,j^*(l)-1)$, os últimos a serem cobertos, satisfazem várias leis exatas
adicionais. Para $W\ge l-1$, escreva
\[
U(l,W) := \Bigl\{\, \textstyle\sum_{i=0}^{l-1} 2^{\beta_i}3^i \bmod 3^l \ \Bigm|\
W\ge\beta_0>\cdots>\beta_{l-1}\ge0 \,\Bigr\},
\]
um análogo de $l$ termos de $R_{j-1,j}$, construído a partir de uma contagem de termos fixa e um
teto de expoente $W$ crescente.

\begin{lemma}[Identidade de escala de largura]
\label{lem:UW}
Para $j\ge l\ge1$,
\begin{equation}
\label{eq:UW}
R_{j-1,j} \equiv 2^{j-l}\,U(l,j+l-1) \pmod{3^l}.
\end{equation}
\end{lemma}

\begin{proof}
Como $3^i\equiv0\pmod{3^l}$ para $i\ge l$, o resíduo de $r=\sum_{i=0}^{j-1}2^{b_i}3^i\in R_{j-1,j}$
módulo $3^l$ depende apenas de seus $l$ expoentes de topo $b_0>\cdots>b_{l-1}$. Uma $l$-tupla
decrescente $b_0>\cdots>b_{l-1}\ge0$ com $b_0\le2j-1$ se estende a uma $j$-tupla admissível
completa em $\{0,\dots,2j-1\}$ (anexando $j-l$ expoentes adicionais, estritamente menores e
distintos) exatamente quando $b_{l-1}\ge j-l$: anexe $b_{l-1}-1,\dots,b_{l-1}-(j-l)$ quando isso
valer, e nenhuma extensão existe quando não valer, já que menos de $j-l$ inteiros não negativos
estão abaixo de $b_{l-1}$. Escrever $c_i := b_i-(j-l)$ transforma $b_{l-1}\ge j-l$ em $c_{l-1}\ge0$
e $b_0\le2j-1$ em $c_0\le2j-1-(j-l)=j+l-1$, e $\sum_i2^{b_i}3^i = 2^{j-l}\sum_i2^{c_i}3^i$. Assim
as tuplas dos $l$ expoentes de topo realizáveis por $R_{j-1,j}$ são exatamente $2^{j-l}$ vezes as
$l$-tuplas decrescentes contadas por $U(l,j+l-1)$, dando a identidade enunciada módulo $3^l$.
\end{proof}

\begin{proposition}[Existência]
\label{prop:exist}
$j^*(l)$ existe para todo $l\ge1$, com $j^*(l)\le2\cdot3^{l-1}-1$.
\end{proposition}

\begin{proof}
$U(l,\cdot)$ é monótona crescente em $W$: alargar a faixa $\{0,\dots,W\}$ só adiciona tuplas de
expoentes admissíveis. Deslocar cada expoente de uma tupla que realiza $u\in U(l,W)$ para cima por
$1$ dá uma $l$-tupla decrescente com expoente de topo no máximo $W+1$ e expoente de base ao menos
$1>0$, ainda admissível, realizando $2u$; logo $2\,U(l,W)\subseteq U(l,W+1)$ para todo $W\ge l-1$.
A tupla $\beta_i=l-1-i$ ($i=0,\dots,l-1$) é admissível em $W=l-1$ e dá
$u_0:=\sum_i2^{l-1-i}3^i\equiv2^{l-1}\pmod3$, uma unidade, de modo que $u_0\in U(l,l-1)$. Iterando
a inclusão, $2^ku_0\in U(l,l-1+k)$ para todo $k\ge0$. Como $2$ tem ordem $2\cdot3^{l-1}$ módulo
$3^l$ (já usado na Seção~\ref{sec:mpg}), $\{2^ku_0 : 0\le k<2\cdot3^{l-1}\}$ é a classe lateral
$u_0\langle2\rangle=(\Z/3^l\Z)^*$: toda unidade módulo $3^l$ aparece como $2^ku_0$ para algum $k$
nessa faixa, logo está em $U(l,l-1+k)\subseteq U(l,W^\dagger)$, $W^\dagger:=l+2\cdot3^{l-1}-2$, por
monotonicidade. Assim $U(l,W^\dagger)=(\Z/3^l\Z)^*$, e o Lema~\ref{lem:UW} em
$j=W^\dagger-l+1=2\cdot3^{l-1}-1\ge l$ dá $R_{j-1,j}\equiv(\Z/3^l\Z)^*\pmod{3^l}$, logo
$j^*(l)\le2\cdot3^{l-1}-1$.
\end{proof}

Essa cota é exponencial, bem acima da taxa conjecturada $\log_43\cdot l$; ela estabelece apenas que
$j^*(l)$ é uma quantidade finita e bem definida, o que o resto desta seção e o argumento de jogo de
payoff médio da Seção~\ref{sec:mpg} então limitam de forma muito mais estreita.

Pelo Lema~\ref{lem:UW}, que precisa ele mesmo de $j\ge l$, o menor $j\ge l$ com
$U(l,j+l-1)=(\Z/3^l\Z)^*$ é o menor $j\ge l$ no qual $R_{j-1,j}$ cobre; como $j^*(l)\ge l$ para
$l=1,\dots,23$ ($l=1$ pela Tabela~\ref{tab:jstar} diretamente, $l=2,\dots,23$ pelo
Lema~\ref{lem:nosmaller}), com igualdade apenas em $l=1$, isso coincide com o próprio $j^*(l)$ em
toda a faixa em que este artigo trabalha, embora nada aqui descarte cobertura em algum $j<l$ ainda
não verificado para um nível ainda não calculado. Escrevendo $W^*(l):=j^*(l)+l-2$ para a largura
um passo antes disso,
\begin{equation}
\label{eq:HviaK}
H(l,j) = 2^{j-l}\,D(l,j+l-1), \qquad D(l,W) := (\Z/3^l\Z)^* \setminus U(l,W).
\end{equation}

\begin{theorem}[Paridade do último resíduo retido]
\label{thm:parity}
Se $j^*(l)\ge l+1$, todo elemento de $H(l,j^*(l)-1)$ é $\equiv1\pmod3$. Isso cobre $l=2,\dots,23$
(Tabela~\ref{tab:jstar}); a hipótese é inaplicável apenas em $l=1$ ($j^*(1)=1$, fora da faixa de
\eqref{eq:HviaK}), o que nada diz sobre se a conclusão em si valeria ali, e voltaria a ser
inaplicável para qualquer $l$ no qual a taxa conjecturada $\log_43<1$ acabe por empurrar $j^*(l)$
abaixo de $l+1$.
\end{theorem}

\begin{proof}
$U(l,\cdot)$ é monótona crescente em $W$, de modo que $D(l,\cdot)$ é monótona decrescente; todo
$x\in D(l,W^*(l))$ é, por minimalidade de $j^*(l)$, representável na largura $W^*(l)+1$ mas não em
$W^*(l)$. Qualquer testemunha para $x$ na largura $W^*(l)+1$ usando apenas expoentes $\le W^*(l)$
já testemunharia $x$nessa largura menor, de modo que toda tal testemunha tem expoente de topo
$\beta_0 = W^*(l)+1$. Sua contribuição ao valor, no ranque $0$, é $2^{\beta_0}\cdot3^0 =
2^{W^*(l)+1}$, logo $x \equiv 2^{W^*(l)+1} \equiv (-1)^{W^*(l)+1} \pmod3$. Por \eqref{eq:HviaK}, que
precisa de $j^*(l)-1\ge l$, elementos de $H(l,j^*(l)-1)$ são $2^{j^*(l)-1-l}x$ para tal $x$, logo
\[
2^{j^*(l)-1-l}x \equiv (-1)^{j^*(l)-1-l}\cdot(-1)^{W^*(l)+1} = (-1)^{2j^*(l)-2} = 1 \pmod3,
\]
usando $W^*(l)+1 = j^*(l)+l-1$ e que $2j^*(l)-2$ é sempre par.
\end{proof}

\begin{proposition}[Quase extinção: contenção direta]
\label{prop:extinctionfwd}
Se $j^*(l)\ge l+2$, $H(l,j^*(l)-1) \subseteq 2\cdot\{x \in H(l,j^*(l)-2) : x\equiv2\!\!\pmod3\}$.
\end{proposition}

\begin{proof}
Escreva $J:=j^*(l)$. O Teorema~\ref{thm:doubling} em $j'=J-2\ge l$ dá $H(l,J-1)\subseteq2H(l,J-2)$,
de modo que todo $y\in H(l,J-1)$ é igual a $2x$ para um (único, já que $2$ é uma unidade)
$x\in H(l,J-2)$. Pelo Teorema~\ref{thm:parity}, $y\equiv1\pmod3$; como $y=2x$, $2x\equiv1\pmod3$,
isto é, $x\equiv2\pmod3$. Assim $x$ está em $\{x\in H(l,J-2):x\equiv2\pmod3\}$ e $y=2x$ está em duas
vezes esse conjunto.
\end{proof}

\begin{empresult}[Bijeção de quase extinção]
\label{er:extinction}
$H(l,j^*(l)-1) = 2\cdot\{x \in H(l,j^*(l)-2) : x\equiv2\!\!\pmod3\}$, exatamente, verificado em
todo $l=2,\dots,9$ contra a tabela exata, sem exceção. A Proposição~\ref{prop:extinctionfwd} prova
a direção $\subseteq$ incondicionalmente (em todo nível com $j^*(l)\ge l+2$, verdadeiro ao longo
desta faixa); apenas a contenção reversa, de que todo $x\equiv2\pmod3$ em $H(l,j^*(l)-2)$ de fato
sobrevive via $2x$, permanece empírica.
\end{empresult}

O mesmo mecanismo de largura que prova o Teorema~\ref{thm:parity} governa a contenção reversa um
passo antes, através do segundo maior expoente $\beta_1$ da testemunha extremal; ao contrário do
argumento do expoente de topo acima, fixar $\beta_1$ precisa da identidade de largura em um passo
do Lema~\ref{lem:width} abaixo. Não levamos esse argumento até uma prova completa dessa direção, e
o registramos como verificado apenas empiricamente.

Uma congruência natural entre níveis, relacionando $H(l+1,j^*(l+1)-1)$ a $H(l,j^*(l)-1)$, não vale.
Escrevendo $\pi_l$ para a redução módulo $3^l$, a congruência candidata
$x\equiv4\,\pi_l(x)\pmod{3^{l+1}}$ é impossível sempre que $\pi_l(x)$ é uma unidade e $l\ge2$:
reduzindo ambos os lados módulo $3^l$ força $\pi_l(x)\equiv4\,\pi_l(x)\pmod{3^l}$, isto é,
$3\,\pi_l(x)\equiv0\pmod{3^l}$. Verificado exaustivamente em $l=2,\dots,6$, $\pi_l(x)$ nunca está
em $H(l,j^*(l)-1)$ para $x\in H(l+1,j^*(l+1)-1)$ em primeiro lugar, de modo que o padrão não tem
sequer uma instância testemunhando-o nesses níveis também. Nenhuma relação entre níveis, entre
$H(l+1,\cdot)$ e $H(l,\cdot)$, é estabelecida por este artigo.

\begin{proposition}[Contenção de duas classes mod $9$]
\label{prop:mod9}
Para $l\ge2$, escrevendo $J:=j^*(l)$, se $J\ge l+1$ então $H(l,J-1) \bmod 9 \subseteq \{4^J,
4^{J+1}\}\pmod9$.
\end{proposition}

\begin{proof}
Como na prova do Teorema~\ref{thm:parity}, toda testemunha de $x\in D(l,W^*(l))$ na largura
$W^*(l)+1$ tem expoente de topo $\beta_0=W^*(l)+1$; módulo $9$, apenas os termos de ranque $0$ e
ranque $1$ sobrevivem (ranque $i\ge2$ contribui um múltiplo de $3^2$), de modo que
$x\equiv2^{\beta_0}+3\cdot2^{\beta_1}\pmod9$ para o segundo maior expoente $\beta_1$ da testemunha.
Por \eqref{eq:HviaK}, elementos de $H(l,J-1)$ são $y=2^{J-1-l}x$, logo
$y\equiv2^{J-1-l+\beta_0}+3\cdot2^{J-1-l+\beta_1}\pmod9$; usando $\beta_0=W^*(l)+1=J+l-1$, o
primeiro termo é $2^{2J-2}=4^{J-1}\bmod9$, fixo. O segundo termo, $3\cdot2^m\pmod9$ para
$m:=J-1-l+\beta_1$, depende apenas da paridade de $m$ (já que $2^m\bmod3$ tem período $2$),
tomando exatamente os dois valores $3$ ou $6\pmod9$ independentemente do valor real de $\beta_1$;
uma verificação direta dá $\{4^{J-1}+3,4^{J-1}+6\}\equiv\{4^J,4^{J+1}\}\pmod9$ para todo $J$. Logo
$y\bmod9\in\{4^J,4^{J+1}\}$.
\end{proof}

\begin{empresult}[Exclusão da classe inferior]
\label{er:mod9}
A classe de fato ocupada é $4^{J+1}$, nunca $4^J$, em todo nível verificado, $l=2,\dots,16$.
Excluir $4^J$ em geral é a questão aberta concreta mais aguda que esta parte da investigação
produziu: a Proposição~\ref{prop:mod9} não impõe restrição alguma sobre o próprio $\beta_1$,
apenas sobre o resultado dependente da paridade, de modo que descartar uma das duas paridades
precisa de um argumento diferente.
\end{empresult}

\begin{proposition}[Suficiência de maxrun limitado]
\label{prop:h014}
Se $\mathrm{maxrun}(H(l,l+1))$ é limitado por uma constante $C$ para todo $l$, então
$j^*(l) \le l + C + 1$.
\end{proposition}

\begin{proof}
Corolário~\ref{cor:bootstrap} em $j=l+1$.
\end{proof}

A recíproca não está estabelecida, e o argumento natural para ela não funciona: lendo o
Corolário~\ref{cor:bootstrap} de trás para frente dá apenas $\mathrm{maxrun}(H(l,l+1)) \ge j^*(l)-l-1$,
uma cota \emph{inferior} sobre maxrun, que não força maxrun a permanecer limitado quando $j^*(l)-l$
não permanece. Provar a recíproca precisaria de uma cota superior sobre maxrun em termos de
$j^*(l)-l$, equivalente a uma afirmação de tightness para o Corolário~\ref{cor:bootstrap} que a
própria questão de redundância de canto desta seção abaixo trata, e que é falsa em sua forma forte,
por testemunha: relembre $1547$ da Seção~\ref{sec:holdouts}.

Uma expectativa ingênua, de que um conjunto aleatório de resíduos retidos na densidade observada
teria cadeias de duplicação muito mais longas do que o teto plano de $3$--$4$ medido em todo nível
$l=5,\dots,22$, não sobrevive ao contato com a densidade real, rapidamente decrescente, de
$H(l,l+1)$. Um modelo de independência ajustado a essa densidade real, descritivo em vez de um
teste estatístico especificado já que nenhuma família distribucional ou convenção de wraparound é
fixada, mantém seus valores preditos dentro da mesma faixa estreita que a sequência observada
ocupa, sem o descompasso flagrante que uma estimativa ingênua de cadeia de duplicação prediria.
Essa faixa é ela mesma estreita, apenas $3$ ou $4$ sobre toda a faixa calculada $l=5,\dots,22$, de
modo que essa comparação principalmente descarta uma superpredição grosseira pelo modelo ingênuo;
ela não estabelece, por si só, que o modelo de independência acompanhe o nível específico em que a
sequência observada sobe de $3$ para $4$ e recua, apenas que seus valores preditos permanecem
dentro dessa mesma faixa estreita. Essa comparação descritiva não diz nada, por si só, sobre se
$\mathrm{maxrun}(H(l,l+1))$ é limitado; a Proposição~\ref{prop:h014} é o único elo estabelecido
entre essa questão e a taxa de $j^*(l)$.

\begin{lemma}[Identidade de largura em um passo]
\label{lem:width}
Para $W\ge l-1$, seja $\mathrm{Corner}(l,W+1)$ o conjunto de valores realizados por uma tupla
testemunha $\beta_0>\cdots>\beta_{l-1}$ com $\beta_0=W+1$ e $\beta_{l-1}=0$ simultaneamente. Então
\[
U(l,W+1) = U(l,W) \cup 2U(l,W) \cup \mathrm{Corner}(l,W+1).
\]
\end{lemma}

\begin{proof}
$\supseteq$: $U(l,W)\subseteq U(l,W+1)$ por monotonicidade, $2U(l,W)\subseteq U(l,W+1)$ pelo
argumento de deslocamento da prova da Proposição~\ref{prop:exist}, e
$\mathrm{Corner}(l,W+1)\subseteq U(l,W+1)$ já que uma testemunha de canto é, em particular, uma
testemunha na largura $W+1$.

$\subseteq$: seja $u\in U(l,W+1)$, testemunhado por $W+1\ge\beta_0>\cdots>\beta_{l-1}\ge0$. Se
$\beta_0\le W$, a tupla já testemunha $u$ na largura $W$, de modo que $u\in U(l,W)$. Se $\beta_0=W+1$
e $\beta_{l-1}\ge1$, subtrair $1$ de cada expoente dá uma tupla decrescente em $[0,W]$ (expoente de
topo $W$, expoente de base $\ge0$) testemunhando $u/2$, de modo que $u\in2U(l,W)$. Se $\beta_0=W+1$
e $\beta_{l-1}=0$, esse é exatamente o caso de canto, de modo que $u\in\mathrm{Corner}(l,W+1)$.
Esses três casos são exaustivos.
\end{proof}

Em toda largura $W$ na faixa $2l+1\le W\le j^*(l)+l-1$, verificado exaustivamente para
$l=3,\dots,13$ (o extremo superior é a maior largura que a tabela exata alcança em cada nível),
$\mathrm{Corner}(l,W+1)$ não contribui nada de novo: o valor de toda testemunha de canto já está
coberto sem ela. Chame isso de propriedade de \emph{redundância de canto} na largura $W$; se ela
vale em toda largura $W\ge2l+1$ em geral, além dos níveis verificados, é aberto. Ela também vale,
verificado diretamente, na única largura menor $W=2l$ em si, para $l=3,4,5,6$; de $l=7$ até $13$
ela é verificada como falhando ali. $W=2l$ é exatamente $j=l+1$, o próprio caso da
Proposição~\ref{prop:h014}; a redundância de canto nessa largura é conhecida, de um jeito ou de
outro, em todo nível $l=3,\dots,13$, valendo apenas para $l=3,4,5,6$. O que permanece aberto é a
própria recíproca da Proposição~\ref{prop:h014}: a falha da redundância de canto em $W=2l$ para
$l=7,\dots,13$ não a resolve em nenhum dos dois sentidos, apenas significa que esse mecanismo em
particular não a prova ali.

\begin{proposition}[Redundância de canto implica tightness]
\label{prop:tight}
Fixe $j$ com $l+2\le j\le j^*(l)$. Se redundância de canto vale em toda largura
$2l+1\le W\le j^*(l)+l-1$, então $j^*(l)=j+\mathrm{maxrun}(H(l,j))$ exatamente.
\end{proposition}

\begin{proof}
Redundância de canto na largura $W$ diz $\mathrm{Corner}(l,W+1)\subseteq U(l,W)\cup2U(l,W)$, de
modo que o Lema~\ref{lem:width} dá $U(l,W+1)=U(l,W)\cup2U(l,W)$. Multiplicação por $2$ é uma
bijeção de $(\Z/3^l\Z)^*$, de modo que comuta com complementação e com interseção; tomando
complementos, isso se torna $D(l,W+1)=D(l,W)\cap2D(l,W)$, e empurrando isso através do
deslocamento em \eqref{eq:HviaK} (ele mesmo uma bijeção, pela mesma razão), a igualdade de nível de
largura em $W=j'+l-1$ se torna uma igualdade de conjuntos de resíduos retidos,
$H(l,j'+1)=2H(l,j')\cap4H(l,j')$, afiando a inclusão do Teorema~\ref{thm:doubling} a uma identidade
sempre que redundância de canto vale nessa largura.

Afirmação: para todo $j'$ com $j\le j'<j^*(l)$, $\mathrm{maxrun}(H(l,j'+1))=\mathrm{maxrun}(H(l,j'))-1$,
dada redundância de canto na largura $j'+l-1$. Como $j'<j^*(l)$, minimalidade de $j^*(l)$ dá
$H(l,j')\ne\emptyset$, de modo que $s:=\mathrm{maxrun}(H(l,j'))\ge1$ ao longo de tudo. A
desigualdade $\le$ é incondicional: se $H(l,j'+1)=\emptyset$ ela lê $0\le s-1$, verdadeiro já que
$s\ge1$; caso contrário, uma cadeia máxima de comprimento $\mathrm{maxrun}(H(l,j'+1))\ge1$ em
$H(l,j'+1)$ se estende, pelo Corolário~\ref{cor:chain}, a uma cadeia um elemento mais longa em
$H(l,j')$. Para $\ge$, fixe uma cadeia máxima $x,2x,\dots,2^{s-1}x\in H(l,j')$. Para cada
$k=2,\dots,s$, tanto $2^{k-1}x$ quanto $2^{k-2}x$ estão nessa cadeia, de modo que a igualdade
$H(l,j'+1)=2H(l,j')\cap4H(l,j')$ coloca $2^kx\in H(l,j'+1)$; isso dá as $s-1$ potências consecutivas
$4x,8x,\dots,2^sx\in H(l,j'+1)$, de modo que $\mathrm{maxrun}(H(l,j'+1))\ge s-1$, provando a
afirmação.

As larguras que isso usa, $j'+l-1$ para $j\le j'<j^*(l)$, estão todas em
$[j+l-1,j^*(l)+l-2]\subseteq[2l+1,j^*(l)+l-1]$ (usando $j\ge l+2$), dentro da hipótese,
independentemente do valor real de $\mathrm{maxrun}(H(l,j))$. Telescopando a afirmação sobre cada
passo de $j'=j$ a $j'=j^*(l)-1$ (um número fixo de passos, $j^*(l)-j$, fixado antes de qualquer
coisa disso) dá $\mathrm{maxrun}(H(l,j^*(l)))=\mathrm{maxrun}(H(l,j))-(j^*(l)-j)$. O lado esquerdo é
$0$ ($H(l,j^*(l))=\emptyset$ por definição), de modo que $\mathrm{maxrun}(H(l,j))=j^*(l)-j$, isto
é, $j^*(l)=j+\mathrm{maxrun}(H(l,j))$.
\end{proof}

Como redundância de canto é verificada em toda largura que a prova acima usa para $l=3,\dots,13$, a
Proposição~\ref{prop:tight} é incondicional nesses níveis: a igualdade do Resultado
Empírico~\ref{er:bootstrap} é um teorema provado ali, para todo $j$ com $l+2\le j\le j^*(l)$. Em
geral, se a hipótese dessa implicação vale é exatamente a questão de redundância de canto acima.

Em $l=3,4,5,6$, onde redundância de canto também vale na largura de fronteira $W=2l$, o mesmo
argumento (começando a telescopagem em $j'=l+1$ em vez de $j'\ge l+2$) prova o caso de fronteira
também: $j^*(l)=(l+1)+\mathrm{maxrun}(H(l,l+1))$ exatamente, verificado diretamente ($l=3$:
$4+2=6$; $l=4$: $5+2=7$; $l=5$: $6+3=9$; $l=6$: $7+3=10$). Essa é a identidade por nível por trás
da própria recíproca da Proposição~\ref{prop:h014}, provada nesses quatro níveis especificamente; a
recíproca em si é uma afirmação sobre todo $l$, e essa rota até ela precisa da identidade em todo
nível, esses quatro incluídos (uma cota superior sobre $\mathrm{maxrun}(H(l,l+1))$ mais fraca que
igualdade exata também bastaria, pela discussão acima, mas nenhuma cota assim é estabelecida
aqui). Em $l=7,\dots,13$, redundância de canto em $W=2l$ é verificada como falhando, de modo que
essa rota não alcança a recíproca ali; além de $l=13$, redundância de canto em $W=2l$ é ela mesma
não verificada. A recíproca permanece aberta em todo $l\ge7$.

\section{Discussão}
\label{sec:discussion}

Duas questões nomeadas resumem o que separa os resultados sobre os últimos resíduos retidos acima
de um teorema mais afiado ali. Excluir a classe inferior na contenção da
Proposição~\ref{prop:mod9} (Resultado Empírico~\ref{er:mod9}) transformaria o padrão de classe
única observado em uma afirmação plenamente provada. Provar redundância de canto em toda largura
$W\ge2l+1$ tornaria, pela Proposição~\ref{prop:tight}, o bootstrap do Corolário~\ref{cor:bootstrap}
tight para todo orçamento $l+2\le j\le j^*(l)$ em geral, não apenas nos níveis $l=3,\dots,13$ onde
redundância de canto já está verificada; isso não resolveria, por si só, a recíproca da
Proposição~\ref{prop:h014} em todo nível, que gira em torno do orçamento de fronteira $j=l+1$
($W=2l$), estabelecido acima apenas em $l=3,4,5,6$, onde redundância de canto também vale ali; em
$l=7,\dots,13$, onde ela é verificada como falhando ali, a recíproca permanece aberta por esse
mecanismo, e além de $l=13$ tanto a propriedade quanto a recíproca são abertas. Nenhuma das duas
questões é, à primeira vista, tão difícil quanto a própria Conjectura da Cobertura Fraca; se
alguma das duas é muito mais fácil não é sabido.

Uma soma exponencial finita construída a partir de tuplas de dígitos estritamente decrescentes e
ponderadas exponencialmente é um objeto natural além deste problema específico de cobertura:
estatísticas $L$ de população finita em frequências e pesos não CLT, teoria de Littlewood--Offord
para sequências de coeficientes estruturadas, e o lado analítico da literatura de Collatz via a
discussão de Tao são todas áreas adjacentes. Achados aqui podem reaparecer lá sob outros nomes. A
rota de espectro completo e cancelamento uniforme que este artigo tenta é inatingível diretamente
(Proposição~\ref{prop:sqrtceiling}); estratégias de cancelamento mais fortes, dependentes de
frequência, permanecem não tentadas. A massa excepcional não se concentra em um conjunto esparso na
única escala acessível verificada ($l=18$, $j=16$, um limiar, apenas frequências primitivas). A
intensidade local, lida em uma profundidade suficientemente fina ($l=12,13,14$), se correlaciona
fortemente com quais resíduos resistem sem determiná-lo completamente, e um diagnóstico de
aleatorização de fase separado se desvia de um nulo que respeita as magnitudes de frequência e a
restrição de anulação-fora-de-unidades que elas satisfazem, no único par verificado,
$(l,m)=(14,16)$. Um nulo adicional, sem carregar informação alguma de fase de frequência
primitiva, corresponde à estatística de razão desse diagnóstico, mas o faz apenas ficando aquém do
máximo absoluto e da raiz quadrada média do vetor real por fatores comparáveis em toda execução;
lido nas unidades absolutas que a razão divide, é esse segundo nulo que se desvia dos dados, por
cerca de mil de seus próprios desvios padrão em energia total, um excesso de magnitude sobre o qual
a estatística de razão é silente. Nenhum dos dois nulos, sozinho ou juntos, estabelece o que decide
a cobertura, ou calibra qualquer um dos desvios a um nível de significância. O resultado posterior
e não relacionado de Tao, de que quase toda órbita de Collatz atinge um valor quase limitado
\cite{tao2019}, trabalha com uma variável aleatória de primeira passagem sobre um modelo
probabilístico diferente do mesmo mapa e não incide diretamente sobre cobertura, mas é a instância
mais clara de um novo tipo de técnica mudando o que conta como tratável nesta área, o tipo de
mudança que uma resolução de qualquer uma das duas questões abertas acima precisaria. Cotas de
densidade de predecessores sobre a própria estrutura de órbitas de $3n+1$, mais desenvolvidas na
abordagem de desigualdades de diferença de Krasikov e Lagarias~\cite{krasikovlagarias2003} (uma
cota inferior de $x^{0.84}$ sobre a contagem de inteiros abaixo de $x$ cuja órbita alcança $1$),
visam a mesma questão ampla que a própria redução de Wirsching de cobertura a densidade positiva de
predecessores, por um mecanismo diferente: desigualdades de diferença em nível de órbita, e não a
estrutura de cobertura de uma família finita $R_{j-1,j}$. Nada nessa literatura atualmente se
transfere para $K(l)$ ou $j^*(l)$ diretamente, mas o alvo compartilhado torna uma transferência
mais plausível do que um resultado não relacionado, e se uma existe é ela mesma uma questão aberta.

A própria Conjectura da Cobertura Fraca, e a taxa assintótica exata de $e(l)$, permanecem abertas.

\section{Disponibilidade de código e dados}
\label{sec:code}

O repositório \url{https://github.com/faculdade/weak-covering-conjecture} contém o código e os
dados por trás de todo cálculo reportado aqui, uma pasta por seção. Ele contém a busca de cobertura
que produziu a Tabela~\ref{tab:jstar}, a comparação de modelos por trás da Tabela~\ref{tab:aic}, o
solucionador do jogo de payoff médio, os cálculos de soma exponencial da Seção~\ref{sec:gapA}, e os
cálculos de resíduos retidos das Seções~\ref{sec:holdouts} e~\ref{sec:further}, junto com os
conjuntos de resíduos retidos armazenados nos orçamentos finais de cada nível $l=5,\dots,21$; o
cálculo de $\mathrm{maxrun}(H(l,l+1))$ por trás da discussão da Proposição~\ref{prop:h014}, que
alcança $l=22$, vive como uma ferramenta diretamente executável em sua própria pasta. Os doze
certificados $(\sigma_k,\rho_k,h_k)$ por trás da Tabela~\ref{tab:mpg} estão armazenados ali com um
verificador que reconfere cada um deles contra a desigualdade na prova do Teorema~\ref{thm:mpg},
sem rodar novamente o solucionador. O README de cada pasta dá o que ela verifica, o comando, a
saída esperada, e o que a execução custa. A maioria das verificações leva segundos; os três maiores
níveis da Tabela~\ref{tab:jstar} não. A versão do repositório referenciada ao longo deste artigo
está arquivada, independentemente de quaisquer mudanças posteriores no ramo padrão do repositório,
em \href{https://doi.org/10.5281/zenodo.21877643}{doi:10.5281/zenodo.21877643} (commit
\texttt{17f1b05}).

\subsection*{Agradecimentos}
Assistência de IA foi usada para revisão textual e tradução.

\renewcommand{\refname}{Refer\^encias}
\begin{thebibliography}{99}

\bibitem{wirsching1998}
G.~J. Wirsching, \emph{The Dynamical System Generated by the $3n+1$ Function}, Lecture Notes in
Mathematics, vol.~1681, Springer, Berlin, 1998.

\bibitem{wirsching2003}
G.~J. Wirsching, \emph{On the problem of positive predecessor density in $3n+1$ dynamics}, Discrete
Contin. Dyn. Syst. \textbf{9} (2003), no.~3, 771--787.

\bibitem{tao2019}
T.~Tao, \emph{Almost all orbits of the Collatz map attain almost bounded values}, Forum Math. Pi
\textbf{10} (2022), Paper No.~e12, 1--56; arXiv:1909.03562.

\bibitem{krasikovlagarias2003}
I.~Krasikov and J.~C. Lagarias, \emph{Bounds for the $3x+1$ problem using difference inequalities},
Acta Arith. \textbf{109} (2003), no.~3, 237--258.

\bibitem{ehrenfeuchtmycielski1979}
A.~Ehrenfeucht and J.~Mycielski, \emph{Positional strategies for mean payoff games}, Internat. J.
Game Theory \textbf{8} (1979), no.~2, 109--113.

\bibitem{meyerovitchyoung2026}
T.~Meyerovitch and A.~Young, \emph{Rationality and computability of the covering radius for sofic
shifts}, arXiv:2603.21449 (2026).

\bibitem{meyerovitchyoung2025nonalt}
T.~Meyerovitch and A.~Young, \emph{Non-alternating mean payoff games}, arXiv:2505.02183 (2025).

\bibitem{tao2011blog}
T.~Tao, \emph{The Collatz conjecture, Littlewood--Offord theory, and powers of $2$ and $3$},
blog post, terrytao.wordpress.com, August 25, 2011.

\end{thebibliography}

\end{document}
